\documentclass[final,3p,times]{elsarticle}

\usepackage{amssymb}
\usepackage{amsmath}
\usepackage{booktabs}
\usepackage{tabularx}
\usepackage{longtable}
\usepackage{multirow}
\usepackage[table,xcdraw]{xcolor}
\usepackage{graphicx}
\usepackage{float}
\usepackage{caption}
\captionsetup{font=small, labelfont=bf}
\usepackage{lineno}
\graphicspath{{figures/}}
\usepackage[hidelinks]{hyperref}
\usepackage{placeins}

\journal{Computers in Biology and Medicine}

\begin{document}

\begin{frontmatter}

\title{Early Prediction of Pediatric Sepsis Using Machine Learning,
Deep Learning, and Explainable AI on the PIC Database: A Comparative
Study Using Phoenix Sepsis Criteria}

\author[nmims]{Tanish Porwal\fnref{equal}}
\author[nmims]{Uzair Teli\corref{cor1}\fnref{equal}}
\author[nmims]{Shruti Sharma}

\fntext[equal]{Tanish Porwal and Uzair Teli contributed equally to this work and share first authorship.}

\cortext[cor1]{Corresponding author. Email address: uzairteli7@gmail.com}

\affiliation[nmims]{organization={Department of Engineering,
             SVKM's NMIMS Indore Campus},
             city={Indore},
             country={India}}

\begin{abstract}
Sepsis remains a leading cause of paediatric intensive care mortality, and existing prediction models are largely trained on adult populations using outdated diagnostic definitions, leaving a gap for paediatric-specific, reproducible work on open data. This study develops a hybrid framework for early prediction of paediatric sepsis on the Pediatric Intensive Care (PIC) database, using the Phoenix Sepsis Criteria as the labelling standard. Using PIC database version 1.1.0 ($n = 3{,}163$ infection-restricted patients), a systematic leakage audit across 306 candidate features removed vasoactive medication indicators and measurement-frequency features carrying label information by construction. The resulting 225-feature dataset was trained using six baseline classifiers, a tuned CatBoost model, three deep learning architectures (FT-Transformer, ResNet, TabNet), a CatBoost--FT-Transformer stacking ensemble (the CBFTTS-SA model), and a CTGAN-based synthetic augmentation pipeline with Iman--Conover correlation correction. Tuned CatBoost achieved AUROC 0.9868 (AUPRC 0.9691, Brier Score 0.0385) at a recommended threshold of 0.365, corresponding to 90\% sensitivity and 96.3\% specificity. The CBFTTS-SA model achieved a statistically indistinguishable AUROC of 0.9851 (DeLong $p = 0.388$) and is presented as the study's forward-looking architecture, since transformer-based models are expected to close the gap with gradient boosting as training data scales; CatBoost remains the stronger, fully interpretable choice at the present cohort size. DeLong tests confirmed CatBoost's significant superiority over five non-CatBoost baselines; FT-Transformer alone was statistically non-inferior ($p = 0.169$) but with substantially lower specificity (0.852 vs 0.963). TreeSHAP analysis identified \texttt{lactate\_max}, \texttt{D-dimer\_max}, and \texttt{platelets\_min} as dominant predictors, mapping onto Phoenix cardiovascular and coagulation domains. Subgroup analysis confirmed robust generalisation across all paediatric age bands (AUROC 0.970--1.000). To the best of our knowledge, this is the first published study combining CatBoost, Phoenix labelling, SHAP explainability, and a CatBoost--Transformer stacking ensemble on the PIC database in one reproducible pipeline.
\end{abstract}

\begin{keyword}
pediatric sepsis \sep machine learning \sep CatBoost \sep
Phoenix Sepsis Criteria \sep SHAP \sep explainable AI \sep
PIC database \sep stacking ensemble \sep synthetic data augmentation
\end{keyword}

\end{frontmatter}

\linenumbers
% =============================================================
\section{Introduction}
\label{sec:literature_review}

Sepsis is a life-threatening condition in which the immune
system's response to infection damages the body's own tissues and
organs rather than resolving the infection; in children it can
progress from suspected infection to multi-organ failure within
hours, and delays in recognition of even an hour have been
associated with measurable increases in mortality  a
relationship first quantified in adult septic shock~\cite{kumar2006}
and corroborated in paediatric critical care analyses~\cite{schlapbach2024}.
Globally, sepsis is one of the most fatal conditions in paediatric
intensive care units (PICUs), with an estimated 25 million
paediatric sepsis cases annually and in-hospital mortality ranging
from roughly 7\% for sepsis to over 33\% for septic shock in
lower-resource settings~\cite{sanchezpinto2024, schlapbach2024}. This scale of
burden, combined with the speed at which children deteriorate,
creates a genuine clinical need for prediction systems that can
flag at-risk patients before organ failure becomes irreversible.

Children are not small adults with respect to sepsis. Their
immune systems, organs, and vital-sign norms are still maturing,
and thresholds that indicate danger in adults may fall within the
normal range for young children~\cite{sanchezpinto2024}. Blood
pressure, respiratory rate, and neurological baselines all differ
by age. The bulk of published sepsis ML models do not account for
this: they are trained on adult ICU data, most commonly
MIMIC-III and MIMIC-IV, and these representational assumptions do
not transfer cleanly to children~\cite{le2019, yang2023}.
Pediatric-specific work on open datasets, with rigorous leakage
controls and a clinically current labelling standard, remains rare.

For years, the diagnostic standard was SIRS (Systemic Inflammatory
Response Syndrome) criteria, adapted for paediatric populations by
Goldstein et al.~\cite{goldstein2005}, which labels sepsis when a
child shows abnormality in two or more of heart rate, respiratory
rate, temperature, and white blood cell count. SIRS is
non-specific: fever from a cold, post-surgical inflammation, and
crying can all trigger its thresholds without any underlying
infection, producing a high false-positive rate that leads to
alert fatigue and unnecessary interventions. PELOD-2 (Pediatric
Logistic Organ Dysfunction score) improved on this for measuring
organ-dysfunction severity in critically ill children~\cite{leteurtre2013},
but it was designed to assess existing dysfunction rather than
predict sepsis onset. The Phoenix Sepsis Score, developed by
Sanchez-Pinto et al.\ and validated through an international
consensus effort by Schlapbach et al., replaces the proxy-based
SIRS approach with a direct measure of organ dysfunction across
four domains  respiratory, cardiovascular, coagulation, and
neurological  validated against over three million patient
encounters~\cite{sanchezpinto2024, schlapbach2024}. A total
Phoenix score of two or above indicates sepsis. This study uses
Phoenix as its labelling standard because it is the only
definition that operationalises actual organ failure rather than
systemic inflammation as a stand-in~\cite{sanchezpinto2024}.

An important methodological constraint must be acknowledged up
front. The Pediatric Intensive Care (PIC) database, the data
source for this study, does not record the fraction of inspired
oxygen (FiO\textsubscript{2}) or mechanical ventilation parameters,
making the respiratory Phoenix component non-computable; the
absence of Glasgow Coma Scale (GCS) documentation similarly renders
the neurological component zero throughout the cohort~\cite{zeng2020}.
Sepsis classification in this study therefore rests on the
cardiovascular and coagulation Phoenix domains alone, a limitation
shared with other machine learning studies applying Phoenix
criteria to datasets with incomplete physiological coverage~\cite{chanci2025}.
This does not undermine the validity of the labelling approach,
given that cardiovascular instability and coagulation dysfunction
are consistently identified as the dominant organ-failure
signatures driving mortality and deterioration in paediatric
sepsis.

Machine learning applied to sepsis prediction has expanded
considerably over the past decade: Fleuren et al.\ catalogued 130
published models across 28 studies in a single systematic review, and Yang
et al.\ examined 23 eligible studies with a narrower methodological
lens, both arriving at the conclusion that predictive performance
has improved with richer feature sets and more sophisticated
algorithms~\cite{fleuren2020, yang2023}. This body of work carries
two structural problems. The first is dataset bias: MIMIC-III and
MIMIC-IV dominate training across the field, and adult ICU
physiology does not transfer cleanly to children, whose
age-dependent vital-sign ranges, immune responses, and organ
dysfunction signatures differ substantially from adult
norms~\cite{yang2023, le2019}. The second is methodological
weakness: Deng et al.\ showed that data leakage, inconsistent
outcome definitions, and weak validation designs appear
persistently across clinical ML studies, overstating the actual
predictive capability of many published models~\cite{deng2022}.
Le et al.\ trained one of the earlier paediatric-specific models
on UCSF PICU data, reporting an AUROC of 0.916 at the time of
onset with prediction up to four hours before clinical sepsis
onset, establishing that early prediction was achievable in
children, though replication across hospital systems has been
limited~\cite{le2019}.

Among architectures applied to structured clinical data, gradient
boosted decision trees have repeatedly come out on top. CatBoost,
XGBoost, and LightGBM handle non-linear feature interactions and
mixed variable types without the preprocessing overhead neural
networks require, and their native class-imbalance handling suits
ICU datasets where positive cases are rare, outperforming deep
learning on moderate-sized tabular datasets
repeatedly~\cite{gorishniy2021, chanci2025}. This preference has
been confirmed empirically in the paediatric sepsis literature.
Chanci et al.\ applied CatBoost with Phoenix Sepsis Criteria
labelling on a US hospital dataset and reported an AUROC of 0.98,
including a SHAP explainability analysis, though the underlying
data remains proprietary and inaccessible for external
validation~\cite{chanci2025}. Han et al.\ compared
five tree-based algorithms -- extra trees, random forest, LightGBM,
XGBoost, and CatBoost -- for paediatric bloodstream infection
prediction, with random forest achieving the best performance
(AUROC 0.874 on the derivation set, 0.762 on external
validation)~\cite{han2024}. Shi et al.\
applied an XGBoost and GRU-based recurrent neural network
framework with SHAP interpretability to predict sepsis in
paediatric patients transitioning from the emergency department to
the PICU, achieving an AUROC of 0.915~\cite{shi2025}.

Two problems recur consistently enough across the sepsis ML
literature to warrant treatment as active design decisions rather
than background noise. The first is class imbalance: in the cohort
used here, sepsis-positive patients are outnumbered roughly three
to one. SMOTE and other synthetic oversampling strategies are the
standard response, but t-SNE visualisation of this dataset's
feature space showed substantial class overlap between
sepsis-positive and sepsis-negative patients, and generating
synthetic minority samples in a region where the classes are not
cleanly separable risks inflating performance metrics without
improving genuine discriminative capability; class weights were
used instead, a standard approach for clinical
ML~\cite{deng2022}. The second is leakage, which Deng et al.\
identified as widespread and underreported across clinical ML
studies~\cite{deng2022}; a systematic pre-training audit of
candidate features was conducted in this study, detailed in
Section~\ref{subsec:leakage}.

Deep learning has been applied to sepsis prediction with
meaningful success, though much of that success has come from
sequential or time-series formulations. Lauritsen et al.\ trained
LSTM-based networks on electronic health record
sequences~\cite{lauritsen2020}; Moor et al.\ applied Gaussian
Process Temporal Convolutional Networks with dynamic time warping
to early sepsis recognition~\cite{moor2019}; Kamaleswaran et al.\
used convolutional neural networks to predict sepsis up to eight
hours before onset in a paediatric cohort~\cite{kamaleswaran2018}.
These results depend on longitudinal data structures that are not
always available or consistently recorded across institutions.

For purely tabular data, two architectures have attracted more
recent attention. Gorishniy et al.\ introduced the Feature
Tokenizer Transformer (FT-Transformer), which treats each tabular
feature as an individual token processed through multi-head
self-attention, outperforming prior deep learning methods across
multiple tabular benchmarks while remaining competitive with
gradient boosted trees~\cite{gorishniy2021}. Arik and Pfister
introduced TabNet, which uses sequential attention to select
features instance by instance, giving it built-in interpretability
without post-hoc methods~\cite{arik2021}. This study benchmarks
both architectures, plus a ResNet baseline, against tuned CatBoost,
and combines CatBoost with the FT-Transformer in a stacking
ensemble  named the \textbf{CBFTTS-SA model} in this study  that
forms the paper's featured architecture.

The PIC training cohort used here contains 2,530 patients, a
sample size that is workable but small enough to raise the
question of whether synthetic augmentation could improve model
robustness. CTGAN, introduced by Xu et al., is a conditional
generative adversarial network for tabular data using
mode-specific normalisation and conditional sampling~\cite{xu2019},
but it is well documented to generate features from marginal
distributions independently, destroying correlation
structure~\cite{xu2019}. In this dataset the problem was severe:
the Spearman correlation between INR and PT in real patients was
0.999, reflecting their shared biological pathway, while vanilla
CTGAN reduced this to 0.016. A three-stage correction pipeline was
developed for this study: CTGAN trained separately per class,
Iman-Conover rank correlation injection~\cite{iman1982} to restore
the real correlation structure, and a zero-inflation correction for
heavy-zero-mass features.

A model that performs well on a held-out test set is not
automatically useful in a clinical environment; trust requires
that the person acting on a prediction can understand why it was
made. A false positive sepsis alert carries real costs
(unnecessary antibiotics, invasive procedures, alert fatigue), and
a missed true positive can be fatal, so asking a clinician to act
on an unexplained probability score is asking them to outsource
judgement to a system they cannot interrogate~\cite{athukorala2026}.
Regulatory frameworks such as the EU AI Act now classify high-risk
medical AI in ways that effectively require interpretability as a
condition of deployment. Among post-hoc explanation methods, SHAP
has become the most widely adopted: Lundberg and Lee grounded
feature attribution in Shapley values from cooperative game
theory~\cite{lundberg2017}, and SHAP is the only additive
attribution method that simultaneously satisfies the local
accuracy, missingness, and consistency axioms~\cite{lundberg2017}.
For tree-based models, the TreeSHAP algorithm computes exact
Shapley values without sampling approximation. Prior work on the
PIC database has already demonstrated that SHAP surfaces clinically
meaningful patterns from Chinese PICU data: Shen et al.\ applied
CatBoost with SHAP for AKI mortality prediction~\cite{shen2025},
Zhou et al.\ used SHAP for DIC prediction~\cite{zhou2025}, and
Huang et al.\ incorporated SHAP into a sepsis survival
analysis~\cite{huang2025}. Outside the PIC context, Yang et al.\
applied explainable ensemble ML with SHAP to MIMIC-IV sepsis data
to identify clinically relevant feature contributions~\cite{jiang2023}.

The PIC database, described by Zeng et al., is the first publicly
accessible intensive care database drawn from a Chinese paediatric
population, maintained by the Children's Hospital of Zhejiang
University School of Medicine~\cite{zeng2020}. Its public
availability makes it one of the few resources outside North
America and Europe on which paediatric clinical ML models can be
openly developed and validated, and it has been used across a
range of prediction tasks: sepsis survival analysis~\cite{huang2025},
AKI mortality with CatBoost and SHAP~\cite{shen2025}, DIC
prediction~\cite{zhou2025}, ensemble-based ICU mortality
prediction~\cite{hong2021}, and the triglyceride-glucose index in
paediatric sepsis mortality~\cite{ding2026}. Cohort construction
and database characteristics for the present study are detailed in
Section~\ref{subsec:dataset_preprocessing}.

Despite this volume of research, a specific and reproducible
combination has not appeared in the published literature. Three
points illustrate the gap most directly:

\begin{itemize}
    \item Chanci et al.\ applied CatBoost with Phoenix Sepsis
    Criteria labelling and reached strong predictive performance
    (AUROC 0.98) with a SHAP explainability layer, but the work
    relied on proprietary US hospital data, so neither the dataset
    nor the trained model is accessible for external
    validation~\cite{chanci2025}.
    \item Shen et al.\ used CatBoost with SHAP on the PIC database,
    but the target outcome was acute kidney injury mortality rather
    than sepsis onset, so the clinical domain and labelling logic
    are entirely different~\cite{shen2025}.
    \item Huang et al.\ worked with PIC data on a sepsis-related
    question but framed it as survival analysis, without Phoenix
    criteria for case definition and without any explainability
    component~\cite{huang2025}.
\end{itemize}

No published study has brought all four elements together in a
single pipeline: CatBoost as a primary model, Phoenix Sepsis
Criteria as the labelling standard, SHAP-based explainability, and
the PIC database as the data source. Each missing element in the
prior work represents a real limitation, whether in
reproducibility, clinical validity of the outcome definition, or
the absence of an interpretive layer clinicians can engage with.

This study fills that gap with a complete machine learning, deep
learning, and explainable AI pipeline for early paediatric sepsis
prediction on the PIC database, using Phoenix Sepsis Criteria as
the labelling standard throughout. To the best of the authors'
knowledge, this is the first published study combining CatBoost,
Phoenix Sepsis Criteria labelling, SHAP explainability, and the PIC
database in a single pipeline; as authors of this work, we developed
a hybrid framework composed of the following four components:

\begin{itemize}
    \item A systematic leakage audit in which 306 candidate
    features were audited before any model was trained, identifying
    and removing features that carried label information by
    construction.
    \item A principled deep learning benchmark in which three
    architectures  FT-Transformer, ResNet, and TabNet  were
    benchmarked against CatBoost using DeLong significance testing,
    and CatBoost was combined with the FT-Transformer in a stacking
    ensemble, named the \textbf{CBFTTS-SA model} in this study and
    presented as the paper's featured architecture, on the basis
    that transformer-class models are expected to close the gap
    with gradient boosting as training data scales beyond this
    cohort's moderate size.
    \item A synthetic data augmentation study in which CTGAN~\cite{xu2019}
    was extended with Iman-Conover rank correlation injection~\cite{iman1982}
    to address CTGAN's known correlation-destruction problem, and
    evaluated honestly against real-data training.
    \item A SHAP interaction analysis in which a lactate--platelets
    axis and an INR--D-dimer coagulopathy axis were identified,
    mapping directly onto the cardiovascular and coagulation
    Phoenix domains and providing biological grounding for the
    CatBoost component's predictions.
    \item A performance comparison in which tuned CatBoost is the
    strongest and fully interpretable model at the present cohort
    size (AUROC 0.9868), while the CBFTTS-SA model is statistically
    indistinguishable from it (AUROC 0.9851, $p = 0.388$) and is
    identified as the architecture with the greatest scaling
    potential, on the basis that transformer-class models are
    expected to close the gap with gradient boosting as training
    data grows beyond this cohort's moderate size  leading to the
    conclusion that CatBoost remains the stronger choice for
    immediate deployment today, but the CBFTTS-SA model is positioned to
    become the superior architecture as larger paediatric ICU
    datasets become available.
\end{itemize}

The remainder of this paper is organised as follows:
Section~\ref{sec:methods} describes the dataset, preprocessing, and
methods; Section~\ref{sec:results} presents experimental results;
Section~\ref{sec:discussion} discusses findings and limitations;
and Section~\ref{sec:conclusion} concludes.

% =============================================================
\section{Materials and Methods}
\label{sec:methods}

\subsection{Data Pipeline Overview}
\label{subsec:dataset_preprocessing}

Figure~\ref{fig:cohort_pipeline} gives a visual overview of the
six-stage cohort-construction and preprocessing pipeline described
in this subsection, from the raw PIC tables through to the
model-ready dataset. The stepwise counts shown in the figure are
under final reconciliation against Table~\ref{tab:cohort_flow} and
the text below; the entry counts, intermediate cohort sizes, and
sepsis-prevalence figures reported in prose throughout this
subsection are the authoritative values for this study.

\begin{figure}[htbp]
    \centering
    \includegraphics[width=\textwidth]{cohort_preprocessing_pipeline}
    \caption{Overview of the cohort-construction and preprocessing
    pipeline, from raw PIC tables through cohort extraction,
    physiological cleaning, feature engineering, Phoenix labelling,
    dataset construction, and the train/test split with median
    imputation, yielding the final model-ready dataset.}
    \label{fig:cohort_pipeline}
\end{figure}

\subsubsection{Data Source}
\label{subsec:datasource}

This study uses the Pediatric Intensive Care (PIC) database
(version~1.1.0), a publicly available, de-identified collection
of electronic health records from paediatric patients admitted to
the intensive care unit at the Children's Hospital of Zhejiang
University School of Medicine in China~\cite{zeng2020}. The
dataset is accessible to credentialed researchers through
PhysioNet and spans routine clinical care collected between 2010
and 2018, covering patient demographics, laboratory measurements,
vital signs, medications, microbiology results, fluid balance,
and clinical symptoms. All timestamps are shifted to preserve
patient privacy while maintaining the relative temporal ordering
of events, making valid time-series analysis possible.

The database contains 12,881 unique patients across 13,941 ICU
stays and is one of the largest publicly available paediatric
intensive care datasets. Its relational structure mirrors that of
the widely benchmarked MIMIC database, adapted for a paediatric
population. Table~\ref{tab:pic_summary} summarises the key
characteristics of the PIC database.

\begin{table}[htbp]
\centering
\caption{Summary characteristics of the PIC database v1.1.0.}
\label{tab:pic_summary}
\small
\begin{tabular}{ll}
\toprule
\textbf{Attribute} & \textbf{Value} \\
\midrule
Dataset             & Pediatric Intensive Care (PIC) \\
Version             & v1.1.0 \\
Total patients      & 12,881 \\
Total ICU stays     & 13,941 \\
Time period         & 2010--2018 \\
Population          & Paediatric (age 0--18) \\
Data type           & EHR (labs, vitals, meds, symptoms) \\
Clinical tables     & 11 (+ 3 reference tables) \\
Source              & Children's Hospital, Zhejiang University \\
Access              & PhysioNet (credentialed) \\
\bottomrule
\end{tabular}
\end{table}

To construct the feature set required for early sepsis onset
prediction, data were extracted from multiple relational tables
within the PIC database and linked using patient and ICU stay
identifiers to form a unified representation of each admission.
The primary clinical tables used were: \texttt{PATIENTS},
\texttt{ADMISSIONS}, and \texttt{ICUSTAYS} for demographics and
admission metadata; \texttt{LABEVENTS} and \texttt{CHARTEVENTS}
for laboratory results and physiological measurements;
\texttt{PRESCRIPTIONS}, \texttt{INPUTEVENTS}, and
\texttt{OUTPUTEVENTS} for treatment and fluid records;
\texttt{MICROBIOLOGYEVENTS} for infection-related data; and
\texttt{EMR\_SYMPTOMS} and \texttt{SURGERY\_INFO} for clinical
context. Three reference tables  \texttt{D\_ITEMS},
\texttt{D\_LABITEMS}, and \texttt{D\_ICD\_DIAGNOSES}  were
used for variable mapping only and did not contribute directly to
the feature matrices.

\subsubsection{Cohort Definition}
\label{subsec:cohort}

The initial pool comprised 13,941 ICU stays from 12,881 unique
paediatric patients. A sequential filtering and restriction
process was applied to define the analysis cohort. Stays where
the patient's age fell outside the range of 0 to 18 years were
excluded ($n = 42$). Repeat ICU admissions were removed, retaining
only each patient's first stay to ensure observation independence
($n = 1{,}060$ excluded). Stays with a total ICU length of stay
under six hours were excluded as they provide insufficient
physiological observation time for meaningful feature extraction
($n = 392$). The resulting intermediate cohort comprised 12,447
ICU stays. A further restriction to patients with clinically
suspected infection was then applied, as the prediction model is
intended for deployment after infection suspicion rather than
across all PICU admissions indiscriminately ($n = 9{,}284$
excluded). The final modelling cohort comprised 3,163 stays.
Table~\ref{tab:cohort_flow} details each step. The majority of the
initial exclusions arose from repeat ICU admissions (7.6\%),
removed to preserve statistical independence of observations; the
restriction to patients with suspected infection (74.6\% of the
intermediate cohort) reflects the intended clinical deployment
context of the model.

\begin{table}[htbp]
\centering
\caption{Cohort construction: stepwise exclusions and restriction
from the PIC database.}
\label{tab:cohort_flow}
\small
\begin{tabular}{lrr}
\toprule
\textbf{Step} & \textbf{Count} & \textbf{\% of Total} \\
\midrule
Initial ICU stays                   & 13,941 & 100.0\% \\
Excluded: age out of range          &     42 &   0.3\% \\
Excluded: non-first ICU stay        &  1,060 &   7.6\% \\
Excluded: LOS $<$ 6 hours           &    392 &   2.8\% \\
\midrule
Intermediate cohort                 & 12,447 &  89.3\% \\
Excluded: no suspected infection    &  9,284 &  66.6\% \\
\midrule
\textbf{Final modelling cohort}     & \textbf{3,163} & \textbf{22.7\%} \\
\bottomrule
\end{tabular}
\end{table}

\subsubsection{Preprocessing Pipeline}
\label{subsec:preprocessing_pipeline}

Data preprocessing was implemented as a six-stage modular
pipeline  cohort filtering, physiological data cleaning and
validation, feature extraction, sepsis label definition, label
verification and dataset construction, and train-test split with
imputation  with each stage producing intermediate output and
logs to enable full reproducibility. This stage corresponds to the
preprocessing block of the overall study pipeline shown in
Figure~\ref{fig:pipeline_overview}.

\paragraph{Physiological data cleaning and validation.} The
objective of this stage was to remove implausible values arising
from measurement or transcription errors while preserving genuine
extreme clinical observations in critically ill patients. For each
variable, physiologically informed upper and lower bounds were
defined with conservative buffers based on reported clinical
extremes in the paediatric critical care
literature~\cite{sanchezpinto2024, zeng2020}; values outside these
bounds were treated as missing rather than corrected, avoiding
distortion of subsequent statistical summaries. Three systematic
data issues were identified and corrected iteratively: haemoglobin
values stored in g/L were being compared against g/dL thresholds,
invalidating 87.9\% of readings, and a unit conversion factor was
applied to resolve this; the anion gap lower bound was revised
from 0 to $-5$~mmol/L to accommodate measurement and calculation
artefacts; and the SpO\textsubscript{2} upper bound was raised
from 100.0\% to 100.5\% to accommodate analyser rounding
artefacts.

\paragraph{Feature extraction.} Feature extraction was performed
over a fixed 24-hour window following ICU admission, chosen
because it aligns with the natural rhythm of clinical assessment
cycles in PICUs and is consistent with the prediction horizon used
in prior paediatric sepsis ML studies~\cite{le2019, chanci2025}.
All predictors are computed from data within this window, ensuring
they temporally precede the prediction target and eliminating any
possibility of future-data leakage. For each physiological
variable, five statistical summaries were computed (maximum,
minimum, mean, first recorded value, measurement count); beyond
these, domain-informed engineering produced trend features
(difference between first and last recorded values), derived
physiological metrics (shock index, mean arterial pressure),
age-adjusted vital-sign abnormality indicators, fluid balance
features, infection-related features derived from antibiotic and
microbiology records, 31 binary symptom indicators, and surgical
context indicators. This stage produced 306 candidate feature columns. Structural
cleaning -- removal of identifier and free-text columns,
zero-variance fields, redundant unit-duplicate columns, Phoenix
component score columns, and deduplication of
\texttt{\_was\_measured} indicators to one column per variable --
reduced this to a 305-feature dataset prior to the leakage audit
described below. The leakage audit (Section~\ref{subsec:leakage})
removed 80 of these features in a single pass, yielding the final
225-feature clean modelling set used throughout this study.

\subsubsection{Leakage Audit}
\label{subsec:leakage}

Prior to model training, a systematic leakage audit was conducted
across the 305-feature structurally-cleaned dataset
(Section~\ref{subsec:preprocessing_pipeline}). Six candidate
leakage suspects were examined. Four were assessed by training an
isolated classifier on each candidate feature group alone and
measuring its stand-alone discriminative power (isolation AUROC):
a group whose isolation AUROC approaches the full-model AUROC
despite representing only a narrow, clinically indirect signal is
evidence that the group encodes label information by construction
rather than genuine physiology.

Measurement frequency features (\texttt{\_count} columns,
$n = 40$) were found to be directly correlated with the Phoenix
labelling process: sepsis-positive patients had systematically
higher measurement counts, driven by clinical urgency rather than
physiology (isolation AUROC 0.878). Vasoactive medication
indicators ($n = 6$ columns) were the most structurally leaky
category: 73.5\% of sepsis patients received vasoactives compared
to 23.7\% of non-sepsis patients, and these variables are
themselves components of the Phoenix cardiovascular domain that
defines the label (isolation AUROC 0.863). Trend features
($n = 16$) and fluid balance variables ($n = 4$) were assessed and
retained: trend features showed a maximum feature-label correlation
of 0.159 and isolation AUROC of 0.848, indicative of genuine
physiological signal, while fluid balance features achieved an
isolation AUROC of only 0.609, confirming no leakage concern.

Two further diagnostic checks were performed. The deduplicated
\texttt{\_was\_measured} missingness indicators ($n = 33$, one
retained per variable after the structural-cleaning deduplication)
showed the same physician-ordering signature as the count columns
(e.g.\ $r = 0.38$ for \texttt{inr\_trend\_was\_measured},
\texttt{fibrinogen\_max\_was\_measured}, and
\texttt{ddimer\_max\_was\_measured} against the sepsis label) and
were removed in full alongside them. A comparison of training-set
imputation medians against full-cohort medians, run as a diagnostic
for outcome-dependent features, flagged \texttt{los\_hours} (ICU
length of stay) with the largest discrepancy (4.165 hours); this
column was identified as an outcome variable unknowable at the time
of prediction and was dropped. A companion patient-overlap check
between training and test partitions could not be directly verified
because patient identifiers were removed during preprocessing;
overlap is instead structurally prevented by the cohort-definition
rule restricting the cohort to each patient's first ICU stay
(Section~\ref{subsec:cohort}).

In total, 80 features across four leaky categories
(\texttt{\_count}, vasoactive medication, \texttt{\_was\_measured},
and \texttt{los\_hours}) were removed, yielding the final clean
feature set of 225 features. A post-cleaning isolation test on this
reduced set achieved AUROC 0.978 and AUPRC 0.950, confirming that
the retained features carry strong legitimate predictive signal.

\begin{figure}[htbp]
    \centering
    \includegraphics[width=\textwidth]{isolation_test_aurocs}
    \caption{Isolation AUROC values of each leakage group.}
    \label{fig:isolation_test_aurocs}
\end{figure}

\subsubsection{Sepsis Label Definition (Phoenix Criteria)}
\label{subsec:labelling}

Sepsis labels were defined using the Phoenix paediatric sepsis
criteria~\cite{sanchezpinto2024}. As established in
Section~\ref{sec:literature_review}, the respiratory and
neurological Phoenix components were non-computable in the PIC
database (FiO\textsubscript{2}, ventilation parameters, and GCS
are absent) and were set to zero; labelling therefore rested on
the cardiovascular and coagulation Phoenix domains exclusively,
consistent with other PIC-based studies facing the same data
constraint~\cite{chanci2025}. Missing values within the computable
domains were treated as contributing zero points, consistent with
the Phoenix framework's convention that absent measurements
contribute no positive score~\cite{sanchezpinto2024}. The final
sepsis label was assigned to any patient meeting both conditions:
clinically suspected infection and a Phoenix score of two or
above.

Applying this definition to the infection-restricted cohort of
3,163 stays identified 864 sepsis cases (27.3\%) and 2,299
patients with suspected infection but no sepsis (72.7\%). For
context, across the full 12,447-stay intermediate cohort, 864
sepsis cases represent 6.9\% prevalence, with 2,299
suspected-infection/no-sepsis patients (18.5\%) and 9,284 patients
with no suspected infection (74.6\%).

\subsubsection{Dataset Construction and Train-Test Split}
\label{subsec:verification}

A dedicated validation stage confirmed the clinical plausibility
of the derived labels. Comparative analysis between
sepsis-positive and sepsis-negative patients showed expected
clinical patterns: higher lactate levels, lower platelet counts,
greater vasoactive medication use, and substantially higher
in-hospital mortality in the sepsis group (18.1\% versus 4.4\%),
consistent with the pathophysiology of sepsis and providing face
validity for the labelling approach.

Two datasets were constructed from the labelled cohort. Option~A
encompassed the full 12,447-stay intermediate cohort ($\{0: 1.0,
1: 13.4\}$ class weights, reflecting 6.9\% prevalence). Option~B
restricted the population to patients with suspected infection,
yielding 3,163 stays ($\{0: 1.0, 1: 2.7\}$ class weights,
reflecting 27.3\% prevalence). \textbf{Option~B was used for all
primary model training, tuning, and evaluation throughout this
study; Option~A is used exclusively for the cohort-definition
sensitivity analysis reported in Section~\ref{subsec:sensitivity},
in which the primary model architecture is retrained from scratch
on Option~A to test whether the study's conclusions are an
artefact of the infection-restriction design choice.} Option~B was
selected as the primary modelling dataset because it reflects the
realistic clinical deployment context: the prediction model would
be applied after infection has been suspected, not indiscriminately
across all PICU admissions.

The Option~B dataset was divided into training (80\%, $n=2{,}530$)
and testing (20\%, $n=633$) sets using stratified sampling to
preserve the 27.3\% sepsis prevalence across both partitions.
Missing values were imputed using median values computed
exclusively on the training set and applied to both partitions,
preventing any leakage of test-set distributional information into
the imputation process. For variables with a missingness rate above 20\%, binary indicator
features (\texttt{\_was\_measured}) were engineered to capture the
clinical signal associated with whether a test was ordered at all,
a threshold consistent with common practice in clinical ML
studies~\cite{deng2022}; after deduplication to one indicator per
variable, these were subsequently identified as leaking the same
physician-ordering signal as the \texttt{\_count} features
(Section~\ref{subsec:leakage}) and were removed from the model
inputs in full -- no \texttt{\_was\_measured} columns are present
in the final 225-feature set. Phoenix score component columns and
residual identifier columns were likewise excluded during the
leakage-audit stage. The leakage audit
(Section~\ref{subsec:leakage}) is therefore the sole source of the
reduction from the 305-feature structurally-cleaned dataset to the
final 225-feature modelling matrix used throughout this study.

\subsubsection{Synthetic Data Augmentation}
\label{subsec:synthetic_intro}

The Option~B training cohort of 2,530 patients is workable but
modest by deep-learning standards, which raises the question of
whether synthetic data could supplement it. To test this, a
CTGAN-based synthetic patient generation pipeline was built and
evaluated as part of this study, extended with an Iman-Conover
rank correlation correction step to address a known limitation of
standard CTGAN on clinically correlated variables (Section~\ref{sec:literature_review}).
The pipeline, its design rationale, and the resulting real-data
versus synthetic-augmented model comparisons are reported in full
in Section~\ref{subsec:synthetic}.

\subsection{Methods}
\label{subsec:methods_all}

\subsubsection{Experimental Setup}
\label{subsec:setup}

All experiments were conducted on the Option~B infection-only
cohort (2,530 training / 633 test samples, 27.3\% sepsis
prevalence), split with stratification to preserve class balance.
Missing values were imputed using training-set medians exclusively,
applied identically to both partitions. A \texttt{StandardScaler}
was applied to numerical features; binary indicators for care unit
type and gender were excluded from scaling. Class imbalance was
addressed using inverse-prevalence class weights
($\{0:1.0, 1:2.7\}$ for Option~B), rather than SMOTE-style
synthesis, for the reasons discussed in Section~\ref{sec:literature_review}.
All model training and evaluation was performed in a conda
environment using Python~3.10 with CatBoost, XGBoost, LightGBM,
scikit-learn, and Optuna. Five-fold stratified cross-validation was
used during training to assess generalisation prior to final
evaluation on the held-out test set. This study was conducted and
reported in accordance with the TRIPOD (Transparent Reporting of a
multivariable prediction model for Individual Prognosis Or
Diagnosis) statement for prediction model development~\cite{collins2015}.

\subsubsection{Baseline Models}
\label{subsec:baseline_methods}

Six classifiers were trained on the 225-feature clean dataset:
CatBoost, XGBoost, LightGBM, Random Forest, Logistic Regression,
and a Multi-Layer Perceptron (MLP). All models used five-fold
stratified cross-validation, and a fixed operating threshold
targeting 85\% sensitivity was applied to enable a fair comparison
of the precision-recall balance across architectures. The
strongest-performing baseline was carried forward for hyperparameter
tuning as the primary single-model architecture.

\subsubsection{CatBoost Tuning and Alternative Pipelines}
\label{subsec:catboost_methods}

CatBoost was selected for tuning based on its superior baseline
performance. One hundred Bayesian optimisation trials were
conducted using Optuna, with 5-fold cross-validated AUPRC as the
objective; AUPRC was preferred over AUROC because it is more
sensitive to classifier performance at the clinically relevant
high-recall region and less susceptible to inflation under class
imbalance. Two alternative single-pipeline architectures were
tested against the resulting tuned CatBoost model. First, a
tree-based stacking ensemble was constructed with tuned CatBoost,
XGBoost, and LightGBM as base learners and a logistic regression
meta-learner: ten clinically motivated interaction features (e.g.\
lactate$_{\max}\times$platelets$_{\min}$, D-dimer$_{\max}\times$INR$_{\max}$)
were engineered and added to the base feature set, out-of-fold
probabilities from five-fold cross-validation trained the
meta-learner, and all base models were retrained on the full
training set before test-set evaluation. Second, a Phoenix-aware
asymmetric loss function was developed in which false negatives are
penalised proportionally to a patient-level severity proxy computed
from normalised z-scores of seven Phoenix component variables
(lactate$_{\max}$, platelets$_{\min}$ (inverted), INR$_{\max}$,
D-dimer$_{\max}$, $|\text{base\_excess}_{\min}|$, PT$_{\text{mean}}$,
PTT$_{\max}$): the per-sample loss for a positive-class patient is
$\text{loss}_i = -w_i\log(p_i)$ with $w_i = \text{base\_weight} +
\alpha\cdot\text{severity}_i$, and negative-class samples retain
standard cross-entropy. Both $\alpha$ and base\_weight were jointly
optimised with CatBoost hyperparameters via Optuna over 50 trials.

The optimal hyperparameter configuration identified by 100 Optuna
trials is presented in Table~\ref{tab:catboost_hyperparams}.

\begin{table}[htbp]
\centering
\caption{Optimal hyperparameters identified by Optuna (100 trials,
AUPRC-optimised). Best CV AUPRC: 0.9580.}
\label{tab:catboost_hyperparams}
\begin{tabular}{lc}
\toprule
\textbf{Hyperparameter} & \textbf{Optimal Value} \\
\midrule
Iterations              & 510   \\
Learning Rate           & 0.0342 \\
Tree Depth              & 7     \\
L2 Leaf Regularisation  & 3.754 \\
Bagging Temperature     & 0.975 \\
Random Strength         & 0.050 \\
Border Count            & 233   \\
Positive Class Weight   & 1.820 \\
\bottomrule
\end{tabular}
\end{table}

\subsubsection{Primary Model Evaluation Protocol}
\label{subsec:eval_protocol}

The tuned CatBoost model was evaluated using five components.
Pairwise DeLong tests assessed whether AUROC differences between
tuned CatBoost and each comparator (baseline classifiers, deep
learning models, alternative pipelines) were statistically
meaningful. Probability calibration was assessed using reliability
diagrams and Platt scaling. Decision Curve Analysis (DCA) evaluated
clinical net benefit across a range of decision thresholds against
\textit{treat-all} and \textit{treat-none} reference strategies. A
full threshold sensitivity analysis characterised the
sensitivity--specificity trade-off across five operating points
spanning 80--95\% sensitivity, to support an evidence-based clinical
deployment recommendation. Finally, subgroup performance was
evaluated separately across four clinically defined paediatric age
bands  neonates (0--28 days), infants (29 days--1 year), children
(1--12 years), and adolescents (12--18 years)  aligned with
established clinical categories in paediatric critical care.

\subsubsection{Deep Learning Benchmark}
\label{subsec:dl_methods}

Three deep learning architectures were benchmarked against tuned
CatBoost: the Feature Tokenizer Transformer (FT-Transformer), which
embeds each tabular feature as an individual token processed
through multi-head self-attention~\cite{gorishniy2021}; a tabular
ResNet, providing a residual feed-forward baseline that isolates
what architectural depth alone contributes without attention; and
TabNet, which applies sequential instance-wise attention masks to
dynamically select a sparse feature subset per patient~\cite{arik2021}.
Each was trained on the same 225-feature leakage-clean dataset with
identical class weights, five-fold stratified cross-validation
during hyperparameter search, and the same 80/20 stratified
train-test split. Hyperparameters for all three architectures were
optimised using Optuna with five-fold cross-validated AUPRC as the
objective, matching the CatBoost tuning protocol. Performance is
reported at the threshold targeting the 90\% sensitivity operating
point for each architecture, with pairwise DeLong tests against
tuned CatBoost.

\subsubsection{CBFTTS-SA Model: CatBoost--FT-Transformer Stacking Ensemble}
\label{subsec:stacking_methods}

This study's featured architecture, named the \textbf{CBFTTS-SA model},
stacks tuned CatBoost and the FT-Transformer as base learners with a
logistic regression meta-learner trained on their out-of-fold
probability outputs from five-fold stratified cross-validation. The
FT-Transformer was selected as the second base learner over ResNet
and TabNet because it produced the highest standalone AUROC among
the three deep learning architectures and its attention-based
representation is architecturally most distinct from CatBoost's
ordered-boosting structure, satisfying the standard diversity
condition for effective stacking. Two configurations were evaluated:
one trained exclusively on the 2,530-patient real-data training set,
and a second trained on the hybrid dataset from the synthetic
augmentation pipeline (Section~\ref{subsec:synthetic_methods}),
combining all real patients with 10,000 CTGAN-generated patients. In
both configurations the meta-learner received only the two
base-learner probability columns as inputs, forcing it to act purely
as a probability combiner; base learners were retrained on the full
training set before final test-set evaluation. As a secondary
experiment, the FT-Transformer's CLS-token attention weight vector
was also extracted and concatenated with the raw 225-feature vector
to train an attention-augmented CatBoost variant (450 input
dimensions), testing whether the transformer's learned inter-feature
weighting could improve CatBoost's split-finding directly rather
than through stacking.

\subsubsection{Synthetic Data Augmentation Pipeline}
\label{subsec:synthetic_methods}

CTGAN, a conditional generative adversarial network for tabular
data~\cite{xu2019}, was used to generate synthetic patients, with a
three-stage correction pipeline developed to address its known
correlation-destruction problem (Section~\ref{sec:literature_review}):
CTGAN was trained separately on sepsis-positive and sepsis-negative
patients to preserve correct within-class marginal distributions,
Iman-Conover rank correlation injection~\cite{iman1982} restored the
real Spearman correlation matrix within each class, and a
zero-inflation correction was applied to variables with heavy zero
mass. The final synthetic dataset comprised 10,000 generated
patients split proportionally by the training cohort's 27.3\% class
prevalence. Three augmentation configurations were evaluated against
the real-data CatBoost baseline: a CatBoost model trained exclusively
on synthetic data; a hybrid CatBoost trained on all 2,530 real
patients (sample weight 4) combined with the 10,000 synthetic
patients (sample weight 1); and the hybrid-data CBFTTS-SA model configuration
described above. All comparisons used DeLong tests against the
real-data tuned CatBoost model, and cross-validation folds used
during Optuna tuning remained unweighted to avoid inflating
optimisation signal from synthetic samples.

\subsubsection{SHAP Explainability Methodology}
\label{subsec:shap_methods}

SHAP explainability analysis was applied to the CatBoost component
using the TreeSHAP algorithm, which computes exact Shapley values
without sampling approximation~\cite{lundberg2017}. SHAP values were
computed for all 633 test-set patients across the full 225 features.
Four analyses were performed: global feature importance (ranked by
mean absolute SHAP value), dependence analysis for the top features,
individual patient explanation via waterfall plots for three
clinically informative cases (a true positive, a false negative, and
a false positive), and pairwise SHAP interaction analysis for the top
ten features using the TreeSHAP interaction API.

\subsubsection{Sensitivity Analysis Design}
\label{subsec:sensitivity_methods}

To test whether the study's conclusions are an artefact of the
Option~B infection-restriction design choice, the tuned CatBoost
architecture  with hyperparameters fixed from the Option~B
optimisation  was retrained from scratch on Option~A (the full
12,447-stay intermediate cohort, regardless of infection suspicion),
with class weights adjusted for its lower sepsis prevalence
($\{0:1.0, 1:13.4\}$ versus Option~B's $\{0:1.0, 1:2.7\}$). This
design isolates the effect of cohort selection from model
architecture and hyperparameter choices: everything except the
training population is held constant between the two configurations.
All primary training, tuning, and evaluation elsewhere in this study
used Option~B exclusively; Option~A is used only for this sensitivity
check, evaluated on the same held-out test set as the primary model.



Figure~\ref{fig:pipeline_overview} summarises the overall
architecture of the study, from the source database through to
explainability analysis. Data flow proceeds from the PIC database
through preprocessing, three parallel modelling tracks, evaluation,
and SHAP explainability analysis, with each of these stages described
in detail in the subsections that follow.

\begin{figure}[!htbp]
    \centering
    \includegraphics[width=0.90\textwidth]{study_pipeline_overview}
    \caption{Complete research pipeline for early paediatric sepsis
    prediction. Data flow proceeds from the PIC database through
    preprocessing, three parallel modelling tracks, evaluation, and
    SHAP explainability analysis. The tuned CatBoost model (amber,
    starred) is the primary single-model architecture selected for
    downstream evaluation.}
    \label{fig:pipeline_overview}
\end{figure}

% =============================================================
\section{Results}
\label{sec:results}

% -- 5.3 --
\subsection{Baseline Model Comparison}
\label{subsec:baseline}

All tree-based ensemble methods attained AUROC above 0.97 on the
225-feature clean dataset; gradient boosting methods substantially
outperformed Random Forest, Logistic Regression, and MLP. CatBoost
produced the strongest performance at baseline (AUROC 0.9863) with
the lowest Brier score (0.0399), as summarised in
Table~\ref{tab:baseline}, and was carried forward for hyperparameter
tuning.

\begin{table}[htbp]
\centering
\caption{Baseline model comparison on the Option~B held-out test set ($n = 633$), at a decision threshold fixed to 85\% sensitivity for cross-architecture comparability. All models were trained with five-fold stratified cross-validation; the metrics reported here are test-set evaluations. Bootstrap 95\% CIs not computed for baseline configurations; see Table~\ref{tab:dl_results}.}
\label{tab:baseline}
\begin{tabular}{lcccccc}
\toprule
\textbf{Model} & \textbf{AUROC} & \textbf{AUPRC} & \textbf{Sensitivity} & \textbf{Specificity} & \textbf{F1} & \textbf{Brier Score} \\
\midrule
CatBoost            & 0.9863 & 0.9684 & 0.855 & 0.980 & 0.881 & 0.0399 \\
LightGBM            & 0.9806 & 0.9570 & 0.855 & 0.963 & 0.876 & 0.0521 \\
XGBoost             & 0.9780 & 0.9512 & 0.867 & 0.963 & 0.882 & 0.0505 \\
Random Forest       & 0.9656 & 0.9222 & 0.855 & 0.928 & 0.836 & 0.0712 \\
Logistic Regression & 0.9152 & 0.8477 & 0.855 & 0.833 & 0.744 & 0.1038 \\
MLP                 & 0.9212 & 0.8635 & 0.855 & 0.863 & 0.771 & 0.1045 \\
\bottomrule
\end{tabular}
\end{table}

The baseline comparison is visualised in
Figure~\ref{fig:roc_pr_baseline}, displaying ROC and
Precision-Recall curves for all six models. The Precision-Recall
curve is informative given the 27.3\% class prevalence. Gradient
boosting methods maintain high precision across the recall range,
while Logistic Regression and MLP show early signs of degradation.

\begin{figure}[htbp]
    \centering
    \includegraphics[width=\textwidth]{phase1_roc_pr_all_models}
    \caption{ROC (left) and Precision-Recall (right) curves for
    all six baseline classifiers on the Option~B test set.}
    \label{fig:roc_pr_baseline}
\end{figure}

% -- 5.4 --
\subsection{Hyperparameter Tuning  CatBoost}
\label{subsec:catboost}

The optimal hyperparameter configuration identified by 100 Optuna
trials is presented in Table~\ref{tab:catboost_hyperparams}
(Section~\ref{subsec:catboost_methods}).

The tuned model achieved AUROC 0.9868 and AUPRC 0.9691 on the
test set. The Optuna optimisation trajectory is shown in
Figure~\ref{fig:optuna_history}.

\begin{figure}[htbp]
    \centering
    \includegraphics[width=\textwidth]{optuna_optimization_history}
    \caption{Optuna trial AUPRC values across 100 trials.
    Convergence is observed by approximately trial 60, with
    subsequent marginal gains.}
    \label{fig:optuna_history}
\end{figure}

\subsubsection{Tree-Based Stacking Ensemble (CatBoost + XGBoost + LightGBM)}
\label{subsubsec:stacking}

This tree-based ensemble is a distinct experiment from the
CatBoost--FT-Transformer CBFTTS-SA model reported in
Section~\ref{subsec:stacking}; it is evaluated here purely as a
same-family (all-gradient-boosting) alternative to single-model
CatBoost. The ensemble achieved AUROC 0.9851 and AUPRC 0.9680,
marginally below tuned CatBoost on both metrics. The DeLong test
confirmed the difference was not statistically significant
($p > 0.05$), and the added pipeline complexity yielded no
measurable benefit over CatBoost alone.

\subsubsection{Phoenix-Aware Custom Loss}
\label{subsubsec:customloss}

Optuna converged to $\alpha = 2.831$ and base\_weight $= 4.467$.
Despite its alignment with the clinical objective, the custom loss
model underperformed significantly: AUROC 0.9674, AUPRC 0.9313, and
Brier Score 0.1005  more than twice that of tuned CatBoost.

\subsubsection{Model Selection}
\label{subsubsec:modelselection}

Neither alternative pipeline improved on tuned CatBoost, which
remains the strongest and most efficient single-model architecture
in this study. A three-way comparison is presented in
Table~\ref{tab:novel_comparison}; the paper's separately evaluated
featured architecture, the CBFTTS-SA model (CatBoost--FT-Transformer
stacking ensemble), is reported in
Section~\ref{subsec:stacking}.

\begin{table}[htbp]
\centering
\caption{Comparison of tuned CatBoost against the two alternative single-family pipelines on the Option~B test set ($n = 633$) at the recommended threshold ($t = 0.365$). Bootstrap 95\% CIs were not computed for these configurations; point estimates are reported.}
\label{tab:novel_comparison}
\begin{tabular}{lcccccc}
\toprule
\textbf{Configuration} & \textbf{AUROC} & \textbf{AUPRC} & \textbf{Sensitivity} & \textbf{Specificity} & \textbf{F1} & \textbf{Brier Score} \\
\midrule
Tuned CatBoost (primary single model) & \textbf{0.9868} & \textbf{0.9691} & \textbf{0.896} & \textbf{0.963} & \textbf{0.899} & \textbf{0.0385} \\
Tree-Based Stacking Ensemble & 0.9851 & 0.9680 & 0.855 & 0.978 & 0.894 & 0.0422 \\
Custom Loss (Phoenix-aware) & 0.9674 & 0.9313 & 0.855 & 0.930 & 0.839 & 0.1005 \\
\bottomrule
\end{tabular}
\end{table}

% -- 5.5 --
\subsection{Comprehensive Evaluation of the Primary Model}
\label{subsec:primary_eval}

\subsubsection{Statistical Comparison: DeLong Tests}

As Table~\ref{tab:delong} shows, the tuned model was statistically
superior to every baseline: LightGBM ($p = 0.017$), XGBoost
($p = 0.005$), and by a larger margin against Random Forest,
Logistic Regression, and MLP (all $p < 0.001$). The one
non-significant result was against the untuned CatBoost baseline
($p = 0.619$), which is expected: hyperparameter tuning sharpened
the model without fundamentally altering its architecture.

\begin{table}[htbp]
    \centering
    \caption{DeLong test results for pairwise AUROC comparisons
    between tuned CatBoost and all baseline classifiers.}
    \label{tab:delong}
    \renewcommand{\arraystretch}{1.3}
    \begin{tabular}{lcccc}
        \toprule
        \textbf{Comparison} & \textbf{AUROC (Tuned)} & \textbf{AUROC (Baseline)} & \textbf{\textit{p}-value} & \textbf{Significant ($p < 0.05$)} \\
        \midrule
        vs CatBoost Baseline    & 0.9868 & 0.9863 & 0.619    & No  \\
        vs LightGBM             & 0.9868 & 0.9806 & 0.017    & Yes \\
        vs XGBoost              & 0.9868 & 0.9780 & 0.005    & Yes \\
        vs Random Forest        & 0.9868 & 0.9656 & $<0.001$ & Yes \\
        vs Logistic Regression  & 0.9868 & 0.9152 & $<0.001$ & Yes \\
        vs MLP                  & 0.9868 & 0.9212 & $<0.001$ & Yes \\
        \bottomrule
    \end{tabular}
\end{table}

\subsubsection{Calibration}

Pre-calibration, the model achieved a Brier Score of $0.0385$,
indicating strong probabilistic calibration. Platt scaling was
applied as a post-hoc step; however, the calibrated model yielded a
marginally higher Brier Score of $0.0422$, a slight degradation
rather than improvement, consistent with the model's decision
boundary already being well-tuned through Optuna. The reliability
diagram (Figure~\ref{fig:calibration}) shows that the model is
well-calibrated in the intermediate probability range ($0.3$--$0.8$),
with minor over-confidence at very high predicted probabilities.
The pre-calibration model is used for all subsequent threshold
analyses.

\begin{figure}[htbp]
    \centering
    \includegraphics[width=0.85\textwidth]{phase3_fig3_calibration.png}
    \caption{Reliability diagram (calibration curve) for the tuned
    CatBoost model before and after Platt scaling. Pre-calibration
    Brier Score: 0.0385; post-calibration: 0.0422. A perfectly
    calibrated model would follow the diagonal dashed line.}
    \label{fig:calibration}
\end{figure}

\subsubsection{Decision Curve Analysis}

The tuned CatBoost model demonstrated positive net benefit over
both the \textit{treat all} and \textit{treat none} reference
strategies across the clinically relevant threshold range of $0.1$
to $0.8$, confirming that deployment of the model would result in
more appropriate interventions than either extreme clinical
strategy (Figure~\ref{fig:dca}).

\begin{figure}[htbp]
    \centering
    \includegraphics[width=0.85\textwidth]{phase3_fig4_dca.png}
    \caption{Decision Curve Analysis for the tuned CatBoost model.
    Net benefit is plotted against the decision threshold
    probability. The model outperforms both the treat-all and
    treat-none baselines across the range of clinically plausible
    thresholds.}
    \label{fig:dca}
\end{figure}
% -- 5.5.4 
\subsubsection{Threshold Analysis and Clinical Operating Point}
\label{subsec:threshold}
 
Table~\ref{tab:threshold} presents five operating points spanning
80\% to 95\% sensitivity, with corresponding classification metrics
at each decision boundary evaluated on the held-out test set.
 
\begin{table}[htbp]
\centering
\caption{Threshold sensitivity analysis for the tuned CatBoost model
on the Option~B test set ($n = 633$). The $\approx$90\% sensitivity operating
point ($t = 0.365$, achieved sensitivity 0.896) is recommended for clinical deployment. Sens.\ =
sensitivity; Spec.\ = specificity.}
\label{tab:threshold}
\small
\begin{tabular}{lcccccc}
\toprule
\textbf{Operating Point} & \textbf{Threshold} & \textbf{Sens.} &
\textbf{Spec.} & \textbf{PPV} & \textbf{NPV} & \textbf{F1} \\
\midrule
80\% Sensitivity              & 0.577 & 0.815 & 0.993 & 0.979 & 0.935 & 0.890 \\
85\% Sensitivity              & 0.447 & 0.861 & 0.980 & 0.943 & 0.949 & 0.900 \\
\textbf{$\approx$90\% (recommended)} & \textbf{0.365} & \textbf{0.896} &
\textbf{0.963} & \textbf{0.901} & \textbf{0.961} & \textbf{0.899} \\
92\% Sensitivity              & 0.276 & 0.919 & 0.943 & 0.859 & 0.969 & 0.888 \\
95\% Sensitivity              & 0.213 & 0.954 & 0.930 & 0.838 & 0.982 & 0.892 \\
\bottomrule
\end{tabular}
\end{table}
 
A threshold of 0.365, targeting 90\% sensitivity (achieved: 0.896), is
recommended as the primary clinical operating point: sensitivity
0.896, specificity 0.963, PPV 0.901, NPV 0.961, F1 0.899. The high
NPV means patients falling below the threshold can be safely
deprioritised with a missed-sepsis rate of only 3.9\%, clinically
tolerable in a PICU context where routine monitoring continues
regardless of model output, while the 96.3\% specificity limits
unnecessary escalations that would otherwise drive alert fatigue.
Across the full threshold range, PPV drops from 0.979 to 0.838 moving
from the 80\% to 95\% sensitivity point (the cost of capturing more
true positives), while F1 stays stable (0.888--0.900) throughout,
confirming the recommended 90\% point sits near the F1 optimum while
preserving workflow-compatible specificity.
 
\begin{figure}[htbp]
    \centering
    \includegraphics[width=\textwidth]{threshold_sensitivity_analysis}
    \caption{Threshold sensitivity curve showing sensitivity,
    specificity, PPV, and F1 score as a function of decision threshold.
    The recommended operating point ($t = 0.365$) is marked with a
    vertical dashed line.}
    \label{fig:threshold_sensitivity}
\end{figure}
 
% -- 5.5.5 
\subsubsection{Subgroup Analysis by Age Band}
\label{subsec:subgroup}
 
Results by age band are summarised in Table~\ref{tab:subgroup}.
 
\begin{table}[htbp]
\centering
\caption{Subgroup performance by age band at the recommended threshold
($t = 0.365$). Prev.\ = sepsis prevalence within subgroup. The
adolescent subgroup ($n = 19$) is small; perfect discrimination should
be interpreted with caution.}
\label{tab:subgroup}
\small
\begin{tabular}{lcccccccc}
\toprule
\textbf{Age Group} & \textbf{N} & \textbf{Sepsis N} &
\textbf{Prev.} & \textbf{AUROC} & \textbf{Sens.} &
\textbf{Spec.} & \textbf{F1} & \textbf{AUPRC} \\
\midrule
Neonate (0--28d)     & 102 & 17 & 16.7\% & 0.970 & 0.824 & 0.965 & 0.824 & 0.915 \\
Infant (29d--1y)     & 223 & 64 & 28.7\% & 0.989 & 0.859 & 0.969 & 0.887 & 0.977 \\
Child (1--12y)       & 289 & 87 & 30.1\% & 0.988 & 0.931 & 0.955 & 0.915 & 0.972 \\
Adolescent (12--18y) &  19 &  5 & 26.3\% & 1.000 & 1.000 & 1.000 & 1.000 & 1.000 \\
\bottomrule
\end{tabular}
\end{table}
 
The model generalised robustly across all subgroups, with AUROC
ranging from 0.970 in neonates to 1.000 in adolescents. The neonate
group had the lowest prevalence (16.7\%) and the most restrained
performance, consistent with the known physiological atypicality of
neonatal sepsis, though 0.970 remains strong relative to any published
benchmark for this age group. The infant and child subgroups together
constitute 80\% of the test cohort (512 of 633 patients) and both
exceeded AUROC 0.987, confirming discriminative performance is not
concentrated in a single age stratum; the adolescent subgroup's
perfect discrimination on only 19 patients carries a wide confidence
interval and does not warrant strong conclusions beyond an absence of
evidence for degraded performance in older children. SHAP feature
rankings were consistent across all age bands  lactate, D-dimer,
and platelets drive predictions throughout, with the neonate group
showing slightly elevated reliance on haematological features.
 
\begin{figure}[htbp]
    \centering
    \includegraphics[width=0.85\textwidth]{phase3_fig5_subgroup_analysis}
    \caption{AUROC and sensitivity by age band with 95\% confidence
    intervals. The neonate group shows the widest uncertainty interval
    attributable to its smaller sample size. All four age bands
    exceed AUROC 0.97.}
    \label{fig:subgroup_auroc}
\end{figure}
 
\begin{figure}[htbp]
    \centering
    \includegraphics[width=\textwidth]{phase3_fig6_summary_dashboard}
    \caption{Comprehensive evaluation dashboard for the tuned CatBoost
    primary model, summarising ROC and Precision-Recall curves,
    calibration diagram, Decision Curve Analysis, confusion matrix,
    and subgroup AUROC results at the recommended threshold
    ($t = 0.365$).}
    \label{fig:summary_dashboard}
\end{figure}

% -- 5.6 --
\subsection{Deep Learning Model Evaluation}
\label{subsec:dl}

Hyperparameter search converged to the configurations detailed in
Tables~\ref{tab:ftt_hyperparams}--\ref{tab:tabnet_hyperparams}.
The FT-Transformer settled on a shallow two-block architecture
with 256 embedding dimensions and eight attention heads, with
relatively low dropout throughout, reflecting the moderate
training set size. The ResNet required eight blocks with
512-dimensional hidden layers and substantially higher second-layer
dropout (0.444), indicating that aggressive regularisation was
needed to prevent overfitting in the deeper feed-forward stack.
TabNet converged to four decision steps with 64-dimensional output
width, a configuration that is common for clinical tabular
datasets of this scale.

\begin{table}[H]
\centering
\caption{Optimal FT-Transformer hyperparameters identified by
Optuna (best 5-fold CV AUPRC: 0.9453).}
\label{tab:ftt_hyperparams}
\small
\begin{tabular}{lc}
\toprule
\textbf{Hyperparameter} & \textbf{Value} \\
\midrule
Transformer blocks ($n\_blocks$)        & 2 \\
Block embedding dimension ($d\_block$)  & 256 \\
Attention heads                         & 8 \\
Attention dropout                       & 0.142 \\
FFN hidden multiplier                   & 2.338 \\
FFN dropout                             & 0.066 \\
Residual dropout                        & 0.013 \\
Learning rate                           & $4.09 \times 10^{-4}$ \\
Weight decay                            & $7.71 \times 10^{-5}$ \\
Batch size                              & 128 \\
\bottomrule
\end{tabular}
\end{table}

\begin{table}[H]
\centering
\caption{Optimal tabular ResNet hyperparameters identified by
Optuna (best 5-fold CV AUPRC: 0.8594).}
\label{tab:resnet_hyperparams}
\small
\begin{tabular}{lc}
\toprule
\textbf{Hyperparameter} & \textbf{Value} \\
\midrule
Residual blocks ($n\_blocks$)           & 8 \\
Main hidden dimension ($d\_main$)       & 512 \\
Inter-block dimension ($d\_hidden$)     & 512 \\
First-layer dropout                     & 0.204 \\
Second-layer dropout                    & 0.444 \\
Learning rate                           & $1.59 \times 10^{-4}$ \\
Weight decay                            & $1.67 \times 10^{-6}$ \\
Batch size                              & 64 \\
\bottomrule
\end{tabular}
\end{table}

\begin{table}[H]
\centering
\caption{Optimal TabNet hyperparameters identified by Optuna
(best 5-fold CV AUPRC: 0.9049).}
\label{tab:tabnet_hyperparams}
\small
\begin{tabular}{lc}
\toprule
\textbf{Hyperparameter} & \textbf{Value} \\
\midrule
Decision step output width ($n\_d$)             & 64 \\
Attention embedding dimension ($n\_a$)          & 16 \\
Decision steps ($n\_steps$)                     & 4 \\
Feature reuse coefficient ($\gamma$)            & 1.082 \\
Sparsity regularisation ($\lambda_{\text{s}}$)  & $2.41 \times 10^{-3}$ \\
Learning rate                                   & $7.89 \times 10^{-3}$ \\
Batch normalisation momentum                    & 0.221 \\
Attention mask type                             & \texttt{entmax} \\
Batch size                                      & 128 \\
Virtual batch size                              & 64 \\
\bottomrule
\end{tabular}
\end{table}

Test-set results for all three architectures, alongside the tuned
CatBoost reference, are presented in Table~\ref{tab:dl_results}.
The FT-Transformer achieved the highest discriminative performance
among the three deep learning models at AUROC 0.9577 (95\% CI:
0.9404--0.9737) and AUPRC 0.9131 (95\% CI: 0.8748--0.9484).
ResNet and TabNet produced near-identical AUROC values of 0.9311
(95\% CI: 0.9046--0.9538) and 0.9306 (95\% CI: 0.9088--0.9497)
respectively, with TabNet showing the weakest AUPRC at 0.8657.
None of the three architectures approached the tuned CatBoost
AUROC of 0.9868.

\begin{table}[H]
\centering
\caption{Deep learning model performance on the Option~B test set
($n = 633$) at the threshold targeting 90\% sensitivity. DeLong
$p$-values are pairwise comparisons against the tuned CatBoost
primary model. 95\% CIs computed by stratified bootstrap
(1000 iterations).}
\label{tab:dl_results}
\small
\begin{tabular}{lcccccc}
\toprule
\textbf{Model} & \textbf{AUROC} & \textbf{AUPRC} & \textbf{Sens.} &
\textbf{Spec.} & \textbf{F1} & \textbf{Brier} \\
\midrule
FT-Transformer        & 0.9577 & 0.9131 & 0.896 & 0.852 & 0.783 & 0.073 \\
ResNet                & 0.9311 & 0.8744 & 0.902 & 0.813 & 0.752 & 0.094 \\
TabNet                & 0.9306 & 0.8657 & 0.908 & 0.770 & 0.720 & 0.108 \\
\midrule
Tuned CatBoost (ref.) & 0.9868 & 0.9691 & 0.896 & 0.963 & 0.899 & 0.0385 \\
\bottomrule
\end{tabular}
\end{table}

Pairwise DeLong tests against the primary model yielded
$p = 0.169$ for the FT-Transformer, $p = 0.018$ for ResNet, and
$p = 0.014$ for TabNet. The FT-Transformer result does not reach
significance at the conventional 0.05 threshold, meaning its
0.029 AUROC gap relative to CatBoost cannot be statistically
distinguished from sampling variation on a test set of 633
patients. The ResNet and TabNet gaps are statistically confirmed;
both architectures are meaningfully inferior to the primary model
in discriminative terms.

A more clinically consequential distinction emerges when
specificity is examined alongside sensitivity. All three deep
learning models achieve similar or higher sensitivity than
CatBoost, but at substantially lower specificity. The
FT-Transformer matches CatBoost's sensitivity of 0.896 exactly,
but specificity falls to 0.852 from CatBoost's 0.963. ResNet and
TabNet push sensitivity to 0.902 and 0.908 respectively, at the
further cost of specificity dropping to 0.813 and 0.770. At a
prevalence of 27.3\%, this specificity gap translates directly
into false-positive burden: at the same sensitivity level, the
FT-Transformer generates approximately four times as many false alerts
per true positive as CatBoost (computed from reported specificities at
matched sensitivity 0.896: FP/TP$_{\text{FT-T}}$ $\approx 0.44$ vs.
FP/TP$_{\text{CatBoost}}$ $\approx 0.11$). In a PICU deployment this matters
operationally, because alert fatigue driven by low precision is
one of the primary mechanisms through which clinical decision
support systems erode clinician trust and compliance over
time~\cite{athukorala2026}. Probability calibration shows the
same gradient: Brier scores range from 0.073 to 0.108 for the
deep learning models compared to 0.0385 for tuned CatBoost,
indicating that all three produce less reliable probability
outputs as a basis for clinical threshold setting.

These results are broadly consistent with the tabular benchmark
findings of Gorishniy et al., who showed that FT-Transformer
becomes competitive with gradient boosting on datasets of
sufficient scale, but that the advantage reverses at moderate
training sizes where CatBoost's implicit regularisation through
ordered boosting provides more stable generalisation than the
data-hungry attention mechanism~\cite{gorishniy2021}. With 2,530
training patients, this cohort sits in that moderate-size regime.
The TabNet outcome is informative in a specific way: instance-wise
feature selection did not produce a performance gain relative to
the simpler ResNet, suggesting that predictive information is
distributed across the feature set broadly rather than concentrated
in a sparse patient-specific subset. This has a clinical
implication of its own, since it means no small panel of features
is sufficient for reliable early sepsis detection uniformly across
all patients in this cohort.

\begin{figure}[htbp]
    \centering
    \includegraphics[width=\textwidth]{phase6_roc_pr_comparison}
    \caption{ROC (left) and Precision-Recall (right) curves for
    the three deep learning architectures compared against tuned
    CatBoost on the Option~B test set. The FT-Transformer most
    closely approaches the CatBoost reference in both curves.
    ResNet and TabNet overlap substantially throughout the
    high-recall region.}
    \label{fig:dl_roc_pr}
\end{figure}

% -- 5.6.1 --
\subsection{CBFTTS-SA Model: CatBoost--FT-Transformer Stacking Ensemble}
\label{subsec:stacking}

The real-data CBFTTS-SA model achieved AUROC 0.9851 (95\% CI: 0.9770--0.9921),
marginally below tuned CatBoost alone (0.9868); the hybrid
(real~+~synthetic) CBFTTS-SA model configuration achieved AUROC 0.9786 (95\%
CI: 0.9683--0.9875). DeLong tests comparing each configuration
against tuned CatBoost yielded $p = 0.388$ and $p = 0.126$
respectively; neither difference is statistically significant at the
0.05 threshold, so the two architectures are reported side by side
in Table~\ref{tab:tu_comparison} and throughout this paper rather
than as a strict ranking. Table~\ref{tab:tu_comparison} reports the
full clinical evaluation for both CBFTTS-SA model configurations at the
threshold targeting 90\% sensitivity for each model, matching the
protocol used for the deep learning benchmark
(Section~\ref{subsec:dl}); SHAP explainability and the
cohort-definition sensitivity check (Section~\ref{subsec:sensitivity})
were performed on the CatBoost component only, since TreeSHAP and
the cohort-definition question both apply most directly to the
model recommended for immediate deployment; the weight analysis
below shows CatBoost carries the large majority of the ensemble's
discriminative signal. To assess whether the hybrid (real~+~synthetic)
CBFTTS-SA configuration is itself viable as a standalone deployment
candidate rather than only as a statistically-tied alternative to
CatBoost, Decision Curve Analysis, a full five-point threshold
sensitivity analysis, and subgroup analysis by age band were
additionally carried out for that configuration at its own
recommended operating threshold; these extended results are reported
in Section~\ref{subsec:hybrid_extended} below.

\begin{table}[htbp]
\centering
\caption{Tuned CatBoost versus the CBFTTS-SA model (CatBoost--FT-Transformer
stacking ensemble) on the Option~B test set ($n = 633$), at the
threshold targeting 90\% sensitivity for each model. 95\% CIs for
AUROC and AUPRC computed by stratified bootstrap (1000 iterations);
see main text.}
\label{tab:tu_comparison}
\small
\begin{tabular}{lccccccc}
\toprule
\textbf{Configuration} & \textbf{AUROC} & \textbf{AUPRC} &
\textbf{Sens.} & \textbf{Spec.} & \textbf{F1} & \textbf{Brier} &
\textbf{DeLong $p$} \\
\midrule
\textbf{Tuned CatBoost (single model)} & \textbf{0.9868} &
\textbf{0.9691} & \textbf{0.896} & \textbf{0.963} & \textbf{0.899} &
\textbf{0.0385} & Reference \\
CBFTTS-SA model (real data) & 0.9851 & 0.9658 & 0.902 & 0.957 & 0.894 &
0.0428 & 0.388 \\
CBFTTS-SA model (hybrid, real+synthetic) & 0.9786 & 0.9534 & 0.908 & 0.935
& 0.872 & 0.0503 & 0.126 \\
\bottomrule
\end{tabular}
\end{table}

A secondary DeLong comparison against the standalone FT-Transformer
(Section~\ref{subsec:dl}) shows the real-data CBFTTS-SA model performs
significantly better than the transformer component alone
($p = 0.0013$); the hybrid CBFTTS-SA model's advantage over the
hybrid-trained FT-Transformer alone is directionally consistent but
not significant ($p = 0.090$). This indicates that stacking
meaningfully improves on the transformer component even though the
ensemble does not surpass CatBoost alone, supporting the choice of
FT-Transformer as the diversity-contributing base learner rather
than confirming stacking as redundant.

The meta-learner weight ratio assigned by logistic regression to
the two base learners is informative: CatBoost received
approximately three to four times the weight of the FT-Transformer
across both configurations, confirming that at the present cohort
size the ensemble's predictive signal is driven predominantly by
the gradient-boosted component. This is consistent with the
benchmark finding in Section~\ref{subsec:dl} that FT-Transformer-class
architectures close the gap with gradient boosting only as
training data scales~\cite{gorishniy2021}; at $n = 2{,}530$ training
patients, CatBoost's ordered-boosting regularisation still has the
edge. The CBFTTS-SA model is nonetheless presented as this study's
featured architecture (Section~\ref{sec:literature_review}), on the
basis that it is statistically indistinguishable from CatBoost
alone today and is architecturally the more promising candidate to
improve further as larger paediatric ICU datasets become available,
whereas tuned CatBoost alone remains the stronger-performing and
fully interpretable choice for immediate deployment at the present
cohort size. The hybrid CBFTTS-SA model configuration produced a slightly
lower AUROC than the real-data configuration, consistent with the
finding in Section~\ref{subsec:synthetic} that synthetic patients
introduce noise the sample weighting scheme does not fully
suppress.

\begin{figure}[htbp]
    \centering
    \includegraphics[width=\textwidth]{phase7_roc_pr}
    \caption{ROC (left) and Precision-Recall (right) curves for
    both CBFTTS-SA model configurations compared against tuned CatBoost
    on the Option~B test set.}
    \label{fig:stacking_roc_pr}
\end{figure}

\subsubsection{Extended Clinical Evaluation of the Hybrid CBFTTS-SA Model}
\label{subsec:hybrid_extended}

The evaluation above establishes that the hybrid CBFTTS-SA model is
statistically indistinguishable from tuned CatBoost at a single
operating point targeting 90\% sensitivity (Table~\ref{tab:tu_comparison}).
To characterise this configuration more completely as a candidate for
future deployment as paediatric ICU datasets scale, the same
threshold-sensitivity, subgroup, and net-benefit analyses applied to
the primary CatBoost model in Sections~\ref{subsec:threshold}
and~\ref{subsec:subgroup} were repeated for the hybrid CBFTTS-SA model
in isolation, using its own recommended operating threshold of
$t = 0.233$ (targeting $\approx$90\% sensitivity, achieved sensitivity
0.9075) rather than borrowing the CatBoost threshold of $t = 0.365$,
since the two models' predicted-probability distributions are not
identical and applying a single model's threshold to another would
misrepresent its true operating characteristics.

\paragraph{Threshold sensitivity.} Table~\ref{tab:hybrid_threshold}
presents the same five operating points spanning 80--95\% sensitivity
used for the primary model. The pattern mirrors the CatBoost result:
specificity trades off against sensitivity from 0.978 at the 80\%
point down to 0.852 at the 95\% point, while F1 remains comparatively
stable (0.813--0.885) across the range, peaking at the 85\% operating
point (F1 = 0.885) rather than exactly at the 90\% point, though the
difference between the two (0.885 vs 0.872) is marginal.

\begin{table}[htbp]
\centering
\caption{Threshold sensitivity analysis for the hybrid (real~+~synthetic)
CBFTTS-SA model on the Option~B test set ($n = 633$). The
$\approx$90\% sensitivity operating point ($t = 0.233$, achieved
sensitivity 0.908) is used as the model's own recommended threshold
throughout this subsection. Sens.\ = sensitivity; Spec.\ = specificity.}
\label{tab:hybrid_threshold}
\small
\begin{tabular}{lcccccc}
\toprule
\textbf{Operating Point} & \textbf{Threshold} & \textbf{Sens.} &
\textbf{Spec.} & \textbf{PPV} & \textbf{NPV} & \textbf{F1} \\
\midrule
80\% Sensitivity              & 0.605 & 0.804 & 0.978 & 0.933 & 0.930 & 0.863 \\
85\% Sensitivity              & 0.402 & 0.867 & 0.965 & 0.904 & 0.951 & 0.885 \\
\textbf{$\approx$90\% (recommended)} & \textbf{0.233} & \textbf{0.908} &
\textbf{0.935} & \textbf{0.840} & \textbf{0.964} & \textbf{0.872} \\
92\% Sensitivity              & 0.184 & 0.925 & 0.926 & 0.825 & 0.970 & 0.872 \\
95\% Sensitivity              & 0.045 & 0.954 & 0.852 & 0.708 & 0.980 & 0.813 \\
\bottomrule
\end{tabular}
\end{table}

\begin{figure}[htbp]
    \centering
    \includegraphics[width=\textwidth]{cbftts_sa_hybrid_threshold_sensitivity}
    \caption{Threshold sensitivity curves for the hybrid CBFTTS-SA
    model showing sensitivity, specificity, PPV, NPV, and F1 score as
    a function of decision threshold. The recommended operating point
    ($t = 0.233$) is marked with a vertical dashed line.}
    \label{fig:hybrid_threshold_sensitivity}
\end{figure}

\paragraph{Subgroup analysis by age band.} Table~\ref{tab:hybrid_subgroup}
reports performance by age band at the model's recommended threshold.
AUROC ranges from 0.961 in neonates to 1.000 in adolescents, closely
tracking the CatBoost subgroup pattern (0.970--1.000,
Table~\ref{tab:subgroup}) though 0.5--0.9 percentage points lower in
every band except the adolescent group, consistent with the small
but consistent AUROC gap between the hybrid CBFTTS-SA model and
CatBoost observed at the overall cohort level
(Table~\ref{tab:tu_comparison}). As with the primary model, the
neonate subgroup shows the most restrained performance and the
smallest prevalence (16.7\%), and the infant and child subgroups
together account for 80\% of the test cohort, confirming that the
hybrid model's discriminative performance is likewise not concentrated
in a single age stratum.

\begin{table}[htbp]
\centering
\caption{Subgroup performance by age band for the hybrid
(real~+~synthetic) CBFTTS-SA model, at its recommended threshold
($t = 0.233$). Prev.\ = sepsis prevalence within subgroup. As with
the primary model, the adolescent subgroup ($n = 19$) is small and
perfect discrimination should be interpreted with caution.}
\label{tab:hybrid_subgroup}
\small
\begin{tabular}{lcccccccc}
\toprule
\textbf{Age Group} & \textbf{N} & \textbf{Sepsis N} &
\textbf{Prev.} & \textbf{AUROC} & \textbf{Sens.} &
\textbf{Spec.} & \textbf{F1} & \textbf{AUPRC} \\
\midrule
Neonate (0--28d)     & 102 & 17 & 16.7\% & 0.961 & 0.765 & 0.965 & 0.788 & 0.871 \\
Infant (29d--1y)     & 223 & 64 & 28.7\% & 0.982 & 0.859 & 0.962 & 0.880 & 0.962 \\
Child (1--12y)       & 289 & 87 & 30.1\% & 0.983 & 0.966 & 0.901 & 0.880 & 0.963 \\
Adolescent (12--18y) &  19 &  5 & 26.3\% & 1.000 & 1.000 & 0.929 & 0.909 & 1.000 \\
\bottomrule
\end{tabular}
\end{table}

\begin{figure}[htbp]
    \centering
    \includegraphics[width=\textwidth]{cbftts_sa_hybrid_subgroup_by_age}
    \caption{AUROC, sensitivity, and specificity by age band for the
    hybrid CBFTTS-SA model at its recommended threshold ($t = 0.233$).
    All four age bands exceed AUROC 0.96, mirroring the age-band
    pattern observed for the primary CatBoost model
    (Figure~\ref{fig:subgroup_auroc}).}
    \label{fig:hybrid_subgroup_auroc}
\end{figure}

\paragraph{Decision Curve Analysis.} The hybrid CBFTTS-SA model
demonstrated positive net benefit over both the \textit{treat all}
and \textit{treat none} reference strategies across the clinically
relevant threshold range of 0.1 to 0.8, the same range used for the
primary CatBoost model's DCA (Section~\ref{subsec:threshold}), and
in fact remained positive across the entire threshold range tested
(0.01--0.95; Figure~\ref{fig:hybrid_dca}), confirming that the hybrid
configuration would, like CatBoost, translate into more appropriate
clinical decisions than either extreme strategy if deployed. Net
benefit declines gradually and stays close to zero from 0.266 at a
threshold of 0.01 down to 0.058 at 0.95, in contrast with the
treat-all strategy, whose net benefit collapses once the decision
threshold exceeds the cohort prevalence of 27.3\% (treat-all net
benefit crosses zero at approximately threshold 0.27 and falls
sharply thereafter, reaching approximately $-13.5$ at threshold
0.95). Figure~\ref{fig:hybrid_dca} plots the full 0.01--0.95 range
on a single axis, so this divergence dominates the vertical scale
and compresses the model curve into a comparatively flat band near
zero; the substantive comparison is nonetheless unambiguous, since
the model curve never crosses the treat-all curve below the
prevalence threshold and remains above the treat-none baseline
(NB $= 0$) throughout, indicating a consistent clinical advantage
across the plausible range of deployment thresholds rather than
only at the single 90\%-sensitivity operating point emphasised
elsewhere in this paper.

\begin{figure}[htbp]
    \centering
    \includegraphics[width=0.85\textwidth]{cbftts_sa_hybrid_dca}
    \caption{Decision Curve Analysis for the hybrid CBFTTS-SA model.
    Net benefit is plotted against the decision threshold probability
    over the full 0.01--0.95 range tested, alongside the treat-all
    and treat-none reference strategies. Because the treat-all curve
    falls sharply below the cohort's 27.3\% prevalence threshold
    (reaching approximately $-13.5$ at threshold 0.95), the axis is
    dominated by that divergence and the model's own curve appears
    as a comparatively flat band close to zero; see main text for the
    corresponding CatBoost comparison plotted on a clinically-windowed
    axis (Figure~\ref{fig:dca}).}
    \label{fig:hybrid_dca}
\end{figure}

Taken together, these extended results indicate that the hybrid
CBFTTS-SA model's near-parity with CatBoost on overall AUROC
(Table~\ref{tab:tu_comparison}) extends to the clinically relevant
threshold-sensitivity, subgroup, and net-benefit profile: the hybrid
configuration is a coherent, self-consistent clinical tool in its own
right rather than only a statistically-tied summary statistic, which
supports its designation as this study's forward-looking architecture
(Section~\ref{sec:literature_review}) even though tuned CatBoost
remains the stronger and more directly interpretable choice for
deployment at the present cohort size.

\subsection{Attention-Augmented CatBoost}
\label{subsec:attention}

The result was a statistically significant degradation in
performance. The attention-augmented model achieved AUROC 0.9657
on the test set, compared to 0.9868 for tuned CatBoost alone. A
pairwise DeLong test confirmed the gap was not attributable to
sampling variation: $p = 0.014$, with the augmented model performing
significantly worse. The Brier score rose from 0.0385 to 0.0978
alongside the discriminative decline, indicating that probability
calibration deteriorated in parallel.

Examining the features assigned the highest attention mass by the
FT-Transformer reveals why augmentation failed rather than helped.
The five most attention-weighted features were
\texttt{calcium\_max}, \texttt{pao2\_max}, care unit designation
(SICU), \texttt{lactate\_max}, and \texttt{platelets\_min}. The
final two are among the top five features by mean absolute SHAP
value in the CatBoost model, and the others are strong physiological
predictors in their own right: the attention weights were not
identifying novel information, but reproducing, in a noisier
continuous form, the same importance ranking CatBoost's split-based
mechanism already encodes. Adding these 225 attention columns
increased feature dimensionality without adding predictive content.

\begin{figure}[htbp]
    \centering
    \includegraphics[width=\textwidth]{phase8_attention_heatmap}
    \caption{FT-Transformer CLS-token attention weight distribution
    across the 225 input features. Attention mass concentrates on
    \texttt{calcium\_max}, \texttt{pao2\_max},
    \texttt{lactate\_max}, and \texttt{platelets\_min}  features
    that are already among the strongest predictors in the CatBoost
    model.}
    \label{fig:attention_heatmap}
\end{figure}

This outcome is a useful negative result: it suggests that the
FT-Transformer and CatBoost, despite their architectural differences,
draw on substantially overlapping predictive structure within this
feature set, so injecting the transformer's summary of that structure
into CatBoost as additional features adds correlated inputs without
a corresponding gain in discriminative signal~\cite{gorishniy2021}.
This is a distinct finding from the CBFTTS-SA model result above: stacking
the two architectures' independent predictions (Section~\ref{subsec:stacking})
is statistically neutral, whereas concatenating the FT-Transformer's
internal representations directly into CatBoost's input space is
actively harmful.

% -- 5.6.4 --
\subsection{Synthetic Data Augmentation}
\label{subsec:synthetic}

As a diagnostic step, vanilla CTGAN was first applied to the full
training dataset before any correction. The result exposed a severe
limitation: the Spearman correlation between INR and PT in real
training data (0.999) collapsed to 0.016 after vanilla CTGAN
generation, meaning a synthetic patient could exhibit a normal INR
alongside a massively prolonged PT  a physiologically impossible
combination. The three-stage correction pipeline
(Section~\ref{subsec:synthetic_methods}) substantially repaired
this correlation structure.

Three augmentation configurations were evaluated against the
real-data tuned CatBoost model (Table~\ref{tab:synthetic_results}).
A CatBoost model trained exclusively on the 10,000 synthetic
patients achieved AUROC 0.9290 ($p < 0.0001$, significantly worse).
A hybrid CatBoost combining all 2,530 real patients (weight 4) with
the synthetic patients (weight 1) improved to AUROC 0.9776
($p = 0.0007$, still significantly worse). The hybrid CBFTTS-SA model
configuration, in which the hybrid dataset trained the base
learners before the meta-learner combined them, reached AUROC
0.9786  the only augmentation configuration whose DeLong test did
not reach significance ($p = 0.126$), because the meta-learner's
dominant real-data CatBoost weighting limits how much synthetic
noise reaches the final output.

\begin{table}[htbp]
\centering
\caption{Synthetic data augmentation results. All DeLong
$p$-values compare against the real-data tuned CatBoost model
(AUROC 0.9868).}
\label{tab:synthetic_results}
\small
\begin{tabular}{lccc}
\toprule
\textbf{Configuration} & \textbf{AUROC} & \textbf{DeLong $p$} &
\textbf{Significant?} \\
\midrule
Synthetic-only CatBoost      & 0.9290 & $< 0.0001$ & Yes (worse) \\
Hybrid CatBoost (weighted)   & 0.9776 & 0.0007     & Yes (worse) \\
$\star$ Hybrid CBFTTS-SA model (Proposed) & 0.9786 & 0.126 & No \\
\midrule
Tuned CatBoost (real data)   & 0.9868 &            & Reference \\
\bottomrule
\end{tabular}
\end{table}

\begin{figure}[htbp]
    \centering
    \includegraphics[width=\textwidth]{hybrid_three_way_roc_pr}
    \caption{Three-way ROC and Precision-Recall comparison between
    the synthetic-only, hybrid-weighted, and real-data CatBoost
    configurations on the Option~B test set.}
    \label{fig:hybrid_roc_pr}
\end{figure}

The clear pattern across all three configurations is monotonic:
AUROC climbs from 0.929 (pure synthetic) to 0.978 (hybrid,
real-weighted) to 0.986 (real data only) as the real-data
contribution to training increases. This is evidence that real
data, not synthetic data, is doing the work  synthetic patients
act as a diluent rather than a booster. The underlying limitation
is not pairwise correlation structure, which the Iman-Conover step
restores reasonably well, but the long-tail distributions of
abnormal laboratory values (e.g.\ \texttt{lactate\_max},
\texttt{D-dimer\_max}) that carry the strongest sepsis signal: a
generative model trained on a few hundred positive examples
systematically underestimates probability mass in the extreme right
tail, and no post-hoc correlation correction can recover information
the training data never provided. The practical implication is that
synthetic augmentation is not a substitute for real data collection
in this setting.

\begin{figure}[htbp]
    \centering
    \includegraphics[width=0.85\textwidth]{distribution_comparison_key_features}
    \caption{Comparison of key feature distributions between real
    training data and CTGAN-generated synthetic data after
    Iman-Conover correlation injection. Marginal distributions are
    broadly preserved, but the extreme right tails of
    \texttt{lactate\_max} and \texttt{D-dimer\_max} are
    underrepresented in the synthetic cohort.}
    \label{fig:synthetic_distribution}
\end{figure}

% -- 5.6.5 --
\subsection{SHAP Explainability Analysis}
\label{subsec:shap}
 
SHAP explainability analysis was applied to the CatBoost component,
which the weight analysis in Section~\ref{subsec:stacking} shows
carries the large majority of the CBFTTS-SA model's discriminative signal.
 
\paragraph{Global feature importance.}
Global feature importance was ranked by mean absolute SHAP value across
all test patients, reflecting each feature's average contribution to
predictions regardless of direction. The top 20 features are presented
in Table~\ref{tab:shap_top20}, visualised as a beeswarm plot in
Figure~\ref{fig:shap_beeswarm}.
 
\begin{table}[htbp]
\centering
\caption{Top 20 features by mean absolute SHAP value, computed across
all 633 test-set patients. Rankings reflect global importance regardless
of prediction direction.}
\label{tab:shap_top20}
\small
\begin{tabular}{clll}
\toprule
\textbf{Rank} & \textbf{Feature} & \textbf{Display Name} &
\textbf{Mean $|\text{SHAP}|$} \\
\midrule
1  & \texttt{lactate\_max}      & Lactate (max)                 & 0.992 \\
2  & \texttt{ddimer\_max}       & D-dimer (max)                 & 0.957 \\
3  & \texttt{platelets\_min}    & Platelets (min)               & 0.776 \\
4  & \texttt{inr\_max}          & INR (max)                     & 0.539 \\
5  & \texttt{calcium\_max}      & Calcium (max)                 & 0.519 \\
6  & \texttt{pt\_max}           & PT (max)                      & 0.259 \\
7  & \texttt{calcium\_mean}     & Calcium (mean)                & 0.222 \\
8  & \texttt{glucose\_max}      & Glucose (max)                 & 0.214 \\
9  & \texttt{pao2\_max}         & PaO\textsubscript{2} (max)    & 0.176 \\
10 & \texttt{ph\_min}           & pH (min)                      & 0.161 \\
11 & \texttt{hematocrit\_min}   & Haematocrit (min)             & 0.146 \\
12 & \texttt{platelets\_first}  & Platelets (first recorded)    & 0.146 \\
13 & \texttt{care\_unit\_SICU}  & Care unit: SICU               & 0.124 \\
14 & \texttt{fibrinogen\_min}   & Fibrinogen (min)              & 0.115 \\
15 & \texttt{care\_unit\_CICU}  & Care unit: CICU               & 0.114 \\
16 & \texttt{ddimer\_mean}      & D-dimer (mean)                & 0.108 \\
17 & \texttt{platelets\_mean}   & Platelets (mean)              & 0.091 \\
18 & \texttt{potassium\_min}    & Potassium (min)               & 0.084 \\
19 & \texttt{lactate\_first}    & Lactate (first recorded)      & 0.075 \\
20 & \texttt{bicarbonate\_min}  & Bicarbonate (min)             & 0.073 \\
\bottomrule
\end{tabular}
\end{table}
 
The top five features  lactate, D-dimer, platelets, INR, and PT 
map directly onto the two computable Phoenix domains: elevated maximum
lactate and depressed minimum platelets correspond to the
cardiovascular domain, while elevated INR, D-dimer, and PT correspond
to the coagulation domain~\cite{sanchezpinto2024, schlapbach2024}.
Calcium (max and mean) at ranks 5 and 7 reflects dysregulated
electrolyte balance in severe systemic infection, and PaO\textsubscript{2},
pH, and bicarbonate in the mid-tier capture respiratory and metabolic
acidosis. Prior SHAP analyses on the PIC database have surfaced
similarly coherent feature hierarchies~\cite{shen2025, zhou2025}.
 
\begin{figure}[htbp]
    \centering
    \includegraphics[width=0.7\textwidth]{phase4_shap_beeswarm}
    \caption{SHAP beeswarm plot for the top 20 features by mean
    absolute SHAP value, computed across all 633 test-set patients.
    Each point represents one patient; colour encodes feature value
    (red = high, blue = low). Features mapping to the Phoenix
    cardiovascular domain (\texttt{lactate\_max},
    \texttt{platelets\_min}) occupy the top ranks, followed by the
    coagulation axis (\texttt{ddimer\_max}, \texttt{inr\_max},
    \texttt{pt\_max}).}
    \label{fig:shap_beeswarm}
\end{figure}
 
\paragraph{SHAP dependence analysis.}
For \texttt{lactate\_max}, the SHAP contribution is approximately zero
below 2~mmol/L and rises sharply above 4~mmol/L, consistent with the
clinical threshold for hyperlactataemia in the Phoenix cardiovascular
domain~\cite{sanchezpinto2024}. D-dimer shows a steep positive
contribution above 2~mg/L (coagulopathy burden), platelets (min) shows
a strongly negative contribution below 100$\times10^9$/L
(thrombocytopaenia), and INR contributions rise sharply above 1.5.
Calcium (max) is non-monotonic: both hypo- and hypercalcaemia are
associated with increased sepsis risk, a U-shaped pattern that cannot
be recovered by linear models and that gives a mechanistic reason for
the AUROC gap between gradient boosting and logistic regression in the
baseline comparison.
 
\begin{figure}[htbp]
    \centering
    \includegraphics[width=\textwidth]{phase4_fig3_dependence_plots}
    \caption{SHAP dependence plots for the top features. The x-axis
    shows raw feature value; the y-axis shows SHAP contribution to
    the sepsis prediction. Colour encodes the value of the
    highest-interacting feature identified by SHAP. The non-monotonic
    calcium pattern and the threshold-like behaviour of lactate and
    D-dimer are visible in the respective panels.}
    \label{fig:shap_dependence}
\end{figure}
 
\paragraph{Individual patient explanations.}
Three cases illustrate the model's behaviour at its margins
(Figure~\ref{fig:shap_waterfall}). The true positive (predicted
probability 0.9997) shows elevated lactate, high D-dimer, and
depressed platelets acting in concert. The false negative
(probability 0.365, at threshold) reflects a borderline case: modest
lactate elevation counterbalanced by near-normal coagulation markers,
characteristic of early or atypical sepsis. The false positive
(probability 0.982; true label non-sepsis) shows laboratory findings
mimicking sepsis, possibly a severe non-infectious inflammatory state.
The model's errors are concentrated at the genuine physiological
boundary of the Phoenix criteria, not in clearly normal or abnormal
presentations.
 
\begin{figure}[htbp]
    \centering
    \includegraphics[width=\textwidth]{phase4_fig4_waterfall_plots}
    \caption{SHAP waterfall plots for three selected test patients:
    true positive (top), false negative (middle), and false positive
    (bottom). Red bars push the prediction toward the sepsis class;
    blue bars push toward non-sepsis. $f(x)$ is the final model
    output and $E[f(x)]$ is the baseline log-odds value.}
    \label{fig:shap_waterfall}
\end{figure}
 
\paragraph{Feature interaction analysis.}
The lactate--platelets interaction is the largest pairwise SHAP
interaction in the model (mean absolute interaction 0.035), spanning
the cardiovascular Phoenix domain: tissue hypoperfusion and
consumptive thrombocytopaenia frequently co-occur in the same
haemodynamic cascade~\cite{sanchezpinto2024}. The INR--D-dimer
coagulopathy axis is second largest (0.021), reflecting converging
evidence of clotting factor depletion and fibrin
degradation~\cite{schlapbach2024}.
 
\begin{figure}[htbp]
    \centering
    \includegraphics[width=0.7\textwidth]{phase4_shap_interaction_heatmap}
    \caption{SHAP interaction value heatmap for the top features.
    The lactate--platelets axis (mean absolute interaction 0.035)
    and the INR--D-dimer coagulopathy axis (0.021) are the two
    dominant pairwise interactions, mapping to the Phoenix
    cardiovascular and coagulation domains respectively.}
    \label{fig:shap_interaction}
\end{figure}
 
\paragraph{SHAP summary and clinical interpretation.}
The dominant features  lactate, D-dimer, platelets, INR, and PT 
collectively represent three of the four Phoenix organ systems:
cardiovascular (lactate, platelets), coagulation (D-dimer, INR, PT),
and metabolic (glucose, pH, bicarbonate). The absence of respiratory
and neurological features from the top rankings is consistent with
the known data limitation: FiO\textsubscript{2} and GCS were
unavailable in PIC v1.1.0 and were set to zero in the Phoenix score
computation~\cite{zeng2020}. The interaction analysis shows the model
has internally learned that simultaneous multi-system compromise is
disproportionately informative, consistent with the additive organ
dysfunction logic underlying the Phoenix score itself, providing a
principled basis for clinician trust in the model's decision
logic~\cite{athukorala2026}. Age does not rank highly among the
top-ranked features, suggesting organ dysfunction markers are
sufficient discriminators across the paediatric age range without
requiring demographic adjustment~\cite{sanchezpinto2024}.
 
\begin{figure}[htbp]
    \centering
    \includegraphics[width=\textwidth]{phase4_fig6_shap_by_class}
    \caption{Mean absolute SHAP contributions stratified by true
    class (sepsis vs non-sepsis), showing which features drive the
    separation between classes most strongly. The coagulation and
    cardiovascular features dominate in both directions.}
    \label{fig:shap_by_class}
\end{figure}
% -- 5.7 --
\subsection{Sensitivity Analysis: Option~A vs Option~B Cohort}
\label{subsec:sensitivity}
 
The Option~A training set comprised 9,957 patients with a sepsis
prevalence of 6.9\%, compared to 2,530 patients at 27.3\% for
Option~B. Performance on the same held-out test set is compared in
Table~\ref{tab:sensitivity}.
 
\begin{table}[htbp]
\centering
\caption{Performance comparison between the primary Option~B
(infection cohort) and sensitivity analysis Option~A (full cohort)
models at the fixed threshold $t = 0.365$. Delta = Option~A minus
Option~B.}
\label{tab:sensitivity}
\small
\begin{tabular}{lcccc}
\toprule
\textbf{Metric} & \textbf{Option~B} & \textbf{Option~A} &
\textbf{Delta} & \textbf{Interpretation} \\
\midrule
AUROC            & 0.9868 & 0.9968 & $+0.010$ & A marginally higher \\
AUPRC            & 0.9691 & 0.9644 & $-0.005$ & B marginally higher \\
Sensitivity      & 0.896  & 0.983  & $+0.087$ & A: fewer missed sepsis \\
Specificity      & 0.963  & 0.968  & $+0.005$ & Near identical \\
PPV              & 0.901  & 0.700  & $-0.201$ & A: lower precision \\
NPV              & 0.961  & 0.999  & $+0.038$ & A: stronger rule-out \\
F1               & 0.899  & 0.817  & $-0.082$ & B: better balanced \\
Brier Score      & 0.0385 & 0.0164 & $-0.022$ & A: lower raw calibration error \\
Training $n$     & 2,530  & 9,957  & $+7,427$ & A much larger \\
Sepsis prevalence & 27.3\% & 6.9\% & $-20.4\%$ & B more balanced \\
\bottomrule
\end{tabular}
\end{table}
 
Option~A achieved a marginally higher AUROC (0.9968 vs 0.9868) and
lower Brier Score (0.0164 vs 0.0385), both attributable to its larger
training set and to the relative ease of discriminating sepsis from the
broader non-infected majority: when 74.6\% of the evaluation population
has no infection suspicion at all, many true negatives are clinically
straightforward and easy to classify correctly, inflating discriminative
and calibration metrics beyond what they would be in the actual
deployment context.

The more operationally meaningful comparison is PPV and F1. Option~A's
PPV of 0.700 means three in ten flagged patients are false positives,
compared to one in ten under Option~B; since a sepsis alert model is
activated after infection has been suspected, not across all admissions
indiscriminately, Option~B's PPV of 0.901 is the figure that would
actually govern clinical workload, and its F1 advantage (0.899 vs 0.817)
confirms more balanced performance at the chosen threshold. These
results confirm that Option~B is the appropriate primary analysis
dataset and that the primary findings are robust rather than artefacts
of the specific cohort definition.
 
\paragraph{SHAP consistency across cohorts.}
The core features driving predictions were consistent across both
cohorts: \texttt{lactate\_max}, \texttt{platelets\_min},
\texttt{ddimer\_max}, \texttt{inr\_max}, and \texttt{glucose\_max} all
appeared in the top ten for both configurations. The primary
difference was elevated importance of antibiotic-count and
microbiology-culture-count features in the Option~A model, where the
broader non-infected population makes treatment-related variables more
discriminative. This reinforces the leakage-aware feature selection
decision made for the primary analysis: in the infection-only cohort
these treatment proxies carry less additional signal once infection is
already established, and including them would risk reintroducing proxy
leakage in a subtler form.
 
\begin{figure}[htbp]
    \centering
    \includegraphics[width=0.7\textwidth]{phase5_fig3_shap_A_vs_B}
    \caption{Side-by-side comparison of top-10 SHAP feature
    importances for the Option~A and Option~B models. The core
    physiological features are consistent across both cohorts;
    treatment-proxy features (antibiotic count, culture count)
    rise in importance in Option~A due to the broader non-infected
    comparison group.}
    \label{fig:sensitivity_shap}
\end{figure}

% =============================================================
\section{Discussion}
\label{sec:discussion}

This study presents the first published combination of CatBoost, Phoenix
Sepsis Criteria labelling, TreeSHAP explainability, and the PIC database in
a single pipeline for early paediatric sepsis prediction, alongside the CBFTTS-SA
model, a CatBoost--FT-Transformer stacking ensemble developed as this
study's forward-looking architecture. The tuned CatBoost model achieved
AUROC 0.9868 and AUPRC 0.9691 on a held-out test set of 633 patients, with
strong calibration (Brier Score 0.0385) and clinically actionable operating
characteristics at the recommended threshold; the CBFTTS-SA model was statistically
indistinguishable from it (AUROC 0.9851, $p = 0.388$). These figures compare
favourably with the existing literature: they are consistent with the 0.98
AUROC reported by Chanci et al.\ using CatBoost and Phoenix labelling on
proprietary US hospital data~\cite{chanci2025}, and exceed the 0.916 AUROC
reported by Le et al.\ in a UCSF PICU cohort using a four-hour prediction
horizon~\cite{le2019}, though differences in labelling standard, prediction
window, and cohort composition complicate direct comparison. Across the
baseline comparison, CatBoost, LightGBM, and XGBoost all exceeded AUROC
0.97, consistent with broader evidence that gradient boosted trees dominate
on moderate-sized clinical tabular data with non-linear feature interactions
and class imbalance~\cite{gorishniy2021, fleuren2020}. The leakage audit is
a further methodological contribution: removing vasoactive medication
indicators, measurement frequency features, deduplicated measurement-
indicator columns, and a length-of-stay outcome proxy eliminated 80
features in total, and
the isolation AUROC of the removed count features alone reached 0.878,
confirming that models retaining such variables would have systematically
overstated predictive capability, a concrete response to the leakage
problem Deng et al.\ identified as pervasive and underreported in clinical
machine learning~\cite{deng2022}.

Benchmarking three deep learning architectures against CatBoost produced a
consistent picture. The FT-Transformer achieved AUROC 0.9577, a gap of 0.029
relative to CatBoost that did not reach statistical significance on this
test set ($p = 0.169$), in line with Gorishniy et al., who showed that
FT-Transformer becomes competitive with gradient boosting only at larger
training scale~\cite{gorishniy2021}; with 2,530 training patients, the
present cohort sits in the moderate-size regime where CatBoost's
ordered-boosting regularisation retains an advantage. The CBFTTS-SA model's result
reinforces this reading: stacking was statistically neutral relative to
CatBoost alone, and the meta-learner assigned three to four times more
weight to CatBoost than to the FT-Transformer, indicating the transformer
component is not yet contributing its full share at this training size.
The CBFTTS-SA model is therefore reported as this study's architecture with the
strongest scaling potential rather than its best current performer, while
tuned CatBoost remains the stronger and more immediately deployable model
today. A more clinically consequential finding is the specificity gap
between CatBoost and the deep learning models: all three matched or
exceeded CatBoost's sensitivity but at substantially lower specificity,
with TabNet as low as 0.770 against CatBoost's 0.963 at comparable
sensitivity, translating to roughly four times as many false alerts per
true positive at this cohort's 27.3\% prevalence. Given that alert fatigue
from low-specificity tools is a recognised mechanism through which clinical
decision support systems lose clinician trust~\cite{athukorala2026}, this
balance of sensitivity and specificity, combined with lower calibration
error and simpler training requirements, makes CatBoost the operationally
preferable choice for this application.

The SHAP analysis indicates that the model's decision logic is grounded in
established pathophysiological markers rather than statistical artefacts.
The top five features, lactate, D-dimer, platelets, INR, and PT, map
directly onto the cardiovascular and coagulation Phoenix domains that
define the sepsis labels in this cohort~\cite{sanchezpinto2024,
schlapbach2024}, and the interaction analysis identifies two mechanistically
coherent axes: a lactate-platelets interaction reflecting simultaneous
circulatory failure and consumptive thrombocytopaenia, and an INR-D-dimer
coagulopathy axis reflecting concurrent clotting factor depletion and
fibrin degradation. The non-monotonic calcium relationship identified in
the dependence analysis, in which both hypocalcaemia and hypercalcaemia are
associated with increased risk, would not be recoverable by logistic
regression and illustrates how gradient boosted trees can surface genuine
clinical signal from raw feature values without explicit engineering. The
synthetic augmentation results are consistent with this reliance on
genuine signal: performance climbed monotonically from synthetic-only
(0.929) through hybrid (0.978) to real-data-only (0.986) training,
confirming that real data, not synthetic data, drives performance, since
the Iman-Conover correction repairs CTGAN's correlation collapse but
cannot manufacture the extreme-tail lactate and D-dimer values that carry
the strongest sepsis signal. The age-stratified analysis further supports
the model's robustness: neonates showed the most restrained but still
strong performance (AUROC 0.970), consistent with the physiological
distinctiveness of the neonatal period, while the infant and child
subgroups, comprising 80\% of the test cohort, both exceeded AUROC 0.987,
indicating that performance is not artificially inflated by concentration
in a single age stratum.



% =============================================================
\section{Conclusion}
\label{sec:conclusion}

This study establishes a complete and reproducible machine learning
pipeline for early paediatric sepsis prediction on the PIC database, using
Phoenix Sepsis Criteria as the labelling standard throughout. A tuned
CatBoost model achieves AUROC 0.9868 and AUPRC 0.9691, the strongest
reported performance for Phoenix-labelled paediatric sepsis prediction on
this database to the authors' knowledge, at a recommended clinical
operating threshold of 0.365 (90\% sensitivity, 96.3\% specificity). A
systematic pre-training leakage audit removed four structurally leaky
feature categories measurement-frequency, vasoactive-medication,
measurement-indicator, and length-of-stay columns from a
305-feature structurally-cleaned candidate pool, offering a
methodological template
for future work on the same database. Benchmarking against three deep
learning architectures shows the FT-Transformer is statistically
non-inferior in discrimination but carries a substantially higher
false-positive burden; the CBFTTS-SA model, a CatBoost--FT-Transformer stacking
ensemble, performs comparably to CatBoost alone ($p = 0.388$) and is
presented as the more promising architecture as paediatric ICU datasets
grow, while CatBoost remains the operationally preferable choice today. A
CTGAN augmentation pipeline with Iman-Conover correlation correction shows
that real data, not synthetic data, drives performance in this setting.
TreeSHAP analysis further grounds the model's decision logic in
established paediatric sepsis pathophysiology, supporting clinician trust
in its outputs.

These contributions address three concrete gaps in the existing
literature: the absence of Phoenix-labelled paediatric sepsis prediction
on reproducible open data, the absence of rigorous leakage controls in
published PIC-based models, and the absence of an interpretable layer
clinicians can engage with in deployment contexts. The findings should be
read alongside the constraints of a retrospective, single-centre design
with a modest events-per-variable ratio and a labelling scheme limited to
two of the four Phoenix domains by the data available. Prospective,
multi-site external validation is accordingly the necessary next step
before any clinical deployment, ideally paired with a time-series
reformulation that evaluates risk continuously rather than at a single
fixed window and, as larger paediatric ICU datasets become available,
re-evaluation of the CBFTTS-SA model's scaling advantage. The methodological
foundation established here, from leakage-audited features through
calibrated, interpretable predictions, provides a platform for that work.

% =============================================================
\section*{Acknowledgements}
The authors thank the PhysioNet platform for providing access to
the PIC database.

% =============================================================
\section*{Conflicts of Interest}
The authors declare no conflicts of interest.

% =============================================================
\section*{Data Availability}
The PIC database is publicly available via PhysioNet at
\url{https://physionet.org}. All preprocessing and modelling code
is available at
\url{https://github.com/tanishqkr/pediatric-sepsis-PIC}.

% =============================================================
\section*{Ethics Statement}
This study used the Paediatric Intensive Care (PIC) database
(version~1.1.0), a de-identified, publicly available dataset
accessible via PhysioNet under credentialed researcher access.
No identifiable patient data were used or accessed. In
accordance with the institutional guidelines of SVKM's NMIMS
Indore Campus and the terms of the PhysioNet Credentialed
Health Data Use Agreement, no additional Institutional Review
Board approval was required for this retrospective analysis of
pre-existing, de-identified data.

% ── Bibliography ──────────────────────────────────────────────
\bibliographystyle{elsarticle-num}
\bibliography{references}

\end{document}